\chapter{Complete Manifolds; Hopf-Rinow and Hadamard Theorems}
\label{chap:dc-complete-manifolds-hopf-rinow-hadamard}

Faithful transcription of do Carmo, \emph{Riemannian Geometry}. Each numbered
statement keeps do Carmo's number, kind and (name); proofs keep his explicit
cross-references ("by \cref{prop:dc-ch3-3-6}", "by \cref{cor:dc-ch3-3-9}", "\cref{thm:dc-ch3-3-7}") so the dependency on the Gauss-lemma
chain is recoverable from the text.

\section{Introduction}

So far we have essentially studied local properties of Riemannian manifolds. One
of the most interesting aspects of differential geometry, however, is the
interplay between the \emph{local} properties (those depending on the behavior of the
manifold near a point) and the \emph{global} properties (those depending on the
behavior of the manifold taken as a whole).

In this chapter we begin the study of the relations between local and global
properties. First we define the natural "habitat" of global properties, namely a
\emph{complete} Riemannian manifold $M$ — a manifold in which the geodesics are
defined for all values of their parameter (cf. \cref{def:dc-ch7-2-2}). Formally this
means that for any $p\in M$, $\exp_p$ is defined on all of $T_pM$; intuitively
that the manifold has no holes or boundaries.

What turns out to be very useful on complete manifolds is the fact (Theorem of
Hopf and Rinow) that given any two points of such a manifold there exists a
\emph{minimizing geodesic} joining them. We prove this in \cref{thm:dc-ch7-2-8}, together with
other facts which imply, for instance, that a compact manifold is complete and
that a closed submanifold of a complete manifold is complete.

As an application of the Theorem of Hopf and Rinow, we prove the theorem of
Hadamard, which states that there is a homeomorphism of a complete, simply
connected manifold of dimension $n$, whose sectional curvature satisfies
$K\le 0$, onto $\mathbb{R}^n$. This is an example of the relation between local
and global properties, in which a local condition ($K\le 0$) together with weak
global restrictions (complete and simply connected) imply a strong global
restriction (homeomorphic to $\mathbb{R}^n$).

From now on, except when explicitly mentioned otherwise, all manifolds will be
supposed \emph{connected}.

The Lean development also adopts the standing convention
$\dim M>0$, expressed by a \texttt{NeZero} instance for
$\operatorname{finrank}_{\mathbb R} E$.  The
connected zero-dimensional case is the trivial one-point case and is not
duplicated in the positive-dimensional analytic development below.

\section{Complete manifolds; Hopf-Rinow Theorem}

When we want to study global properties of a differentiable manifold $M$, we need
to ensure that $M$ is not a proper open submanifold of a manifold $M'$. The usual
condition guaranteeing this non-extendibility is compactness; in certain cases,
however, we would like a weaker condition.

\begin{definition}[extendible]
  \label{def:dc-ch7-2-1}
  \dcref{ch7:2.1}
  \lean{Riemannian.IsExtendible}
  \leanok
  A Riemannian manifold $M$ is \emph{extendible} if there exists a Riemannian manifold
  $M'$ such that $M$ is isometric to a proper open subset of $M'$. In the opposite
  case, $M$ is \emph{non-extendible}.

  In Lean this is encoded on the same model space: there is a connected manifold $M'$, a
  metric $g'$ on $M'$, and a smooth injective local diffeomorphism $\varphi:M\to M'$ with
  $\operatorname{range}\varphi\ne M'$ that preserves the metric ($|d\varphi\,v|=|v|$). A
  local diffeomorphism is an open map, so $\operatorname{range}\varphi$ is an open subset,
  and injectivity makes $\varphi$ a diffeomorphism onto it — i.e. an isometry onto a proper
  open subset. Non-extendibility is the negation.
  
  It happens that the class of non-extendible Riemannian manifolds is too large. A
  convenient subfamily is given by the next definition.
\end{definition}

\begin{definition}[geodesically complete]
  \label{def:dc-ch7-2-2}
  \dcref{ch7:2.2}
  \uses{def:dc-ch3-expmap}
  \lean{Riemannian.Exponential.expDomainIntrinsic, Riemannian.Geodesic.expDomainIntrinsic_eq_univ_iff_isGeodesicallyCompleteAt, Riemannian.Geodesic.IsGeodesicallyComplete}
  \leanok
  A Riemannian manifold $M$ is \emph{(geodesically) complete} if for all $p\in M$ the
  exponential map $\exp_p$ is defined for all $v\in T_pM$; i.e. if any geodesic
  $\gamma(t)$ starting from $p$ is defined for all values of the parameter
  $t\in\mathbb{R}$.

  Formally, \texttt{expDomainIntrinsic g p} is the time-one domain of the
  chart-independent maximal geodesic obtained by gluing continuous
  moving-chart geodesics.  The theorem
  \texttt{expDomainIntrinsic\_eq\_univ\_iff\_isGeodesicallyCompleteAt}
  proves, without assuming metric completeness, that this domain is all of
  $T_pM$ exactly when every initial velocity at $p$ has a continuous global
  geodesic.  The older chart-fixed \texttt{expDomain} is only a local flow
  domain and is not used in this definition.
\end{definition}

\begin{lemma}[a local isometry maps geodesics to geodesics]
  \label{lem:dc-ch7-2-3-geodesic-push}
  \dcref{ch7:2.3}
  \uses{def:dc-ch1-2-2, lem:dc-ch7-3-4-rays-are-geodesics}
  \lean{Riemannian.RiemannianMetric.eq_of_metricInner_eq,
    Riemannian.Geodesic.solvesGeodesicODEAt_comp_of,
    Riemannian.Geodesic.isGeodesic_comp_of_isLocalDiffeomorph}
  \leanok
  Let $f:M\to M'$ be a smooth local diffeomorphism that preserves the metric
  ($|df\,v|=|v|$ for all $v$; equivalently $f^*g'=g$). Then $f$ carries geodesics of $g$ to
  geodesics of $g'$: for every continuous $g$-geodesic $\gamma$, the image $f\circ\gamma$ is a
  $g'$-geodesic. This is the geodesic tool behind do Carmo's proof of Proposition 2.3 (``the
  geodesic $\gamma(t)=\tilde\gamma(1-t)$ of $M'$ is a geodesic in $M$''); it is the \emph{push}
  companion of the pullback geodesic transfer of \cref{lem:dc-ch7-3-4-rays-are-geodesics}.
\end{lemma}
\begin{proof}
  \leanok
  \uses{lem:dc-ch7-3-4-rays-are-geodesics}
  Metric preservation identifies $g$ with the pulled-back metric $f^*g'$ — a Riemannian
  metric is determined by its inner product at every point (extensionality,
  \texttt{eq\_of\_metricInner\_eq}) — so each $g$-geodesic solves the $f^*g'$-geodesic ODE.
  Under a local diffeomorphism the Christoffel contraction transforms by the differential
  $df$ alone (the transformation law of \cref{lem:dc-ch7-3-4-rays-are-geodesics}), so applying
  $df$ to the vanishing geodesic operator $u''+\Gamma^{f^*g'}(u',u')$ shows the reading of
  $f\circ\gamma$ satisfies $W''+\Gamma^{g'}(W',W')=df(0)=0$ at every time; no injectivity of
  $df$ is needed for this direction. $\square$
\end{proof}

\begin{proposition}[Proposition 2.3]
  \label{prop:dc-ch7-2-3}
  \dcref{ch7:2.3}
  \uses{def:dc-ch7-2-1, def:dc-ch7-2-2, lem:dc-ch7-2-3-geodesic-push, thm:dc-ch3-3-7}
  \lean{Riemannian.not_isExtendible_of_isGeodesicallyComplete}
  \leanok
  If $M$ is complete then $M$ is non-extendible.
\end{proposition}
\begin{proof}
  \leanok
  \uses{def:dc-ch7-2-1, def:dc-ch7-2-2, lem:dc-ch7-2-3-geodesic-push, thm:dc-ch3-3-7}
  Suppose $M\subset M'$ is isometric to a proper open subset of a
  Riemannian manifold $M'$. Since $M'$ is connected, the boundary $\partial M$ of
  $M$ in $M'$ is non-empty. Let $p\in\partial M$ and let $U'\subset M'$ be a normal
  neighborhood of $p$ in $M'$. Let $q\in U'\cap M$ and let $\tilde\gamma(t)$ be a
  geodesic in $M'$ with $\tilde\gamma(0)=p$, $\tilde\gamma(1)=q$. Then
  $\gamma(t)=\tilde\gamma(1-t)$, $|t| < \delta$, is a geodesic in $M$ with
  $\gamma(0)=q$. This geodesic is not defined for some $t\le 1$, which contradicts
  the completeness of $M$. $\square$
  
  It can be shown by an example that the converse is \emph{not} true (cf. Exercise 4),
  so the class of non-extendible manifolds is actually larger than the class of
  complete manifolds.
  
  At this stage it is convenient to introduce a distance function on a Riemannian
  manifold (not necessarily complete) $M$. Given $p,q\in M$, consider all the
  piecewise differentiable curves joining $p$ to $q$. Since $M$ is connected such
  curves exist (cover a continuous curve joining $p$ to $q$ by finitely many
  coordinate neighborhoods and replace each "piece" contained in a coordinate
  neighborhood by a differentiable curve).
\end{proof}

\begin{definition}[distance]
  \label{def:dc-ch7-2-4}
  \dcref{ch7:2.4}
  \uses{def:dc-ch1-2-9}
  \lean{Manifold.riemannianEDist, Riemannian.RiemannianMetric.IsRiemannianDist}
  The \emph{distance} $d(p,q)$ is defined by $d(p,q)=$ infimum of the lengths of all
  curves $f_{p,q}$, where $f_{p,q}$ is a piecewise differentiable curve joining $p$
  to $q$.

  Mathlib's \texttt{Manifold.riemannianEDist} takes the infimum over $C^1$
  paths.  The formal development uses that $C^1$ presentation throughout;
  the equivalence with do Carmo's larger piecewise differentiable curve class
  has not been formalized, so this node is deliberately not marked
  \texttt{\textbackslash leanok}.
\end{definition}

\begin{lemma}[finiteness of the distance]
  \label{lem:dc-ch7-2-4-finite}
  \dcref{ch7:2.4}
  \uses{def:dc-ch7-2-4}
  \lean{Manifold.riemannianEDist_ne_top}
  \leanok
  If $M$ is connected, then $d(p,q)<\infty$ for all $p,q\in M$.
\end{lemma}
\begin{proof}
  \leanok
  Fix $p$ and let $A=\{q\in M;\ d(p,q)<\infty\}$. Every point $x\in M$ has a
  neighborhood $V_x$ all of whose points are at finite distance from $x$: in a
  chart at $x$, the straight segment from the image of $x$ to the image of a
  nearby point $y$ pulls back to a differentiable curve in $M$ from $x$ to $y$
  whose length is bounded by a constant multiple of the chart distance, the
  differential of the chart inverse being bounded near $x$. If $q\in A$ then
  $V_q\subset A$ by the triangle inequality, and if $q\notin A$ then
  $V_q\cap A=\varnothing$, again by the triangle inequality. Hence $A$ is open
  and closed; since $p\in A$ and $M$ is connected, $A=M$. $\square$
\end{proof}

\begin{lemma}[short curves stay near their origin]
  \label{lem:dc-ch7-2-5-chart}
  \dcref{ch7:2.5}
  \uses{def:dc-ch7-2-4}
  \lean{setOf_riemannianEDist_lt_subset_nhds}
  \mathlibok
  Let $p\in M$ and let $S$ be a neighborhood of $p$. Then there exists $c>0$
  such that every $q\in M$ with $d(p,q)<c$ lies in $S$. (In a chart at $p$,
  fix a closed neighborhood $u\subset S$ of $p$ on which the differential of
  the chart is bounded by $C$, an open $v\subset u$, and $r>0$ with the chart
  ball of radius $r$ contained in the chart image of $v$; a curve of length
  $<r/C$ starting at $p$ has, while it remains in $u$, chart image of length
  $<r$, so it remains in $v$, strictly inside $u$; a continuity induction
  along the curve then keeps it in $u$ for its whole duration, so its
  endpoint lies in $u\subset S$.)
\end{lemma}

\begin{lemma}[positivity of the distance]
  \label{lem:dc-ch7-2-5-pos}
  \dcref{ch7:2.5}
  \uses{def:dc-ch7-2-4}
  \lean{Manifold.riemannianEDist_pos}
  \leanok
  If $p\ne q$ then $d(p,q)>0$.
\end{lemma}
\begin{proof}
  \leanok
  \uses{lem:dc-ch7-2-5-chart}
  Since $M$ is Hausdorff, $S=M\smallsetminus\{q\}$ is a neighborhood of $p$.
  By \cref{lem:dc-ch7-2-5-chart} there is $c>0$ such that every point at
  distance $<c$ from $p$ lies in $S$. Since $q\notin S$, it follows that
  $d(p,q)\ge c>0$. $\square$
\end{proof}

\begin{proposition}[$d$ is a metric]
  \label{prop:dc-ch7-2-5}
  \dcref{ch7:2.5}
  \uses{def:dc-ch7-2-4}
  \lean{Manifold.riemannianEDist_self, Manifold.riemannianEDist_comm, Manifold.riemannianEDist_triangle, MetricSpace.ofRiemannianMetric}
  \leanok
  With the distance $d$, $M$ is a metric space, that is:
  1. $d(p,r)\le d(p,q)+d(q,r)$,
  2. $d(p,q)=d(q,p)$,
  3. $d(p,q)\ge 0$, and $d(p,q)=0\iff p=q$.
\end{proposition}
\begin{proof}
  \leanok
  \uses{lem:dc-ch7-2-4-finite, lem:dc-ch7-2-5-pos}
  (1), (2) and two of the assertions of (3) are immediate consequences of the
  definition of the infimum: reversing a curve preserves its length, the
  concatenation of a curve from $p$ to $q$ with a curve from $q$ to $r$ is a
  piecewise differentiable curve from $p$ to $r$ whose length is the sum of
  the two lengths, the constant curve gives $d(p,p)=0$, and $d$ is finite by
  \cref{lem:dc-ch7-2-4-finite}. It remains to show that if $p\ne q$ then
  $d(p,q)>0$, which is \cref{lem:dc-ch7-2-5-pos}. (do Carmo instead takes a
  normal ball $B_r(p)$ not containing $q$ and invokes the minimizing property
  of radial geodesics, \cref{prop:dc-ch3-3-6}; the chart-comparison proof of
  \cref{lem:dc-ch7-2-5-pos} reaches the same conclusion without normal
  balls.) $\square$

  Observe that if there exists a minimizing geodesic $\gamma$ joining $p$ to $q$
  (which is not always true), then $d(p,q)=\text{length of }\gamma$.
\end{proof}

\begin{proposition}[topology coincides]
  \label{prop:dc-ch7-2-6}
  \dcref{ch7:2.6}
  \uses{def:dc-ch7-2-4}
  \lean{eventually_riemannianEDist_lt, setOf_riemannianEDist_lt_subset_nhds, PseudoEMetricSpace.ofRiemannianMetric}
  \mathlibok
  The topology induced by $d$ on $M$ coincides with the original topology on $M$.
\end{proposition}
\begin{proof}
  \uses{lem:dc-ch7-2-4-finite, lem:dc-ch7-2-5-chart}
  Every metric ball $\{q;\ d(p,q)<c\}$, $c>0$, is a neighborhood of $p$ in the
  manifold topology: as in the proof of \cref{lem:dc-ch7-2-4-finite}, every
  point $q$ of a sufficiently small chart ball around $p$ satisfies
  $d(p,q)<c$, the pulled-back chart segment from $p$ to $q$ having length
  bounded by a constant multiple of the chart distance. Conversely, every
  manifold neighborhood of $p$ contains a metric ball around $p$ by
  \cref{lem:dc-ch7-2-5-chart}. Hence the two topologies have the same
  neighborhoods at every point. $\square$
\end{proof}

\begin{corollary}[distance is continuous]
  \label{cor:dc-ch7-2-7}
  \dcref{ch7:2.7}
  \uses{def:dc-ch7-2-4}
  \lean{Manifold.continuous_riemannianEDist_right}
  \leanok
  If $p_0\in M$, the function $f:M\to\mathbb{R}$ given by $f(p)=d(p,p_0)$ is
  continuous.

  The fact which makes the concept of completeness relevant is the following
  theorem.
\end{corollary}
\begin{proof}
  \leanok
  \uses{prop:dc-ch7-2-5, prop:dc-ch7-2-6}
  By the triangle inequality, $|d(p,p_0)-d(p',p_0)|\le d(p,p')$, so $f$ is
  Lipschitz, hence continuous for the metric topology, which coincides with
  the topology of $M$ by \cref{prop:dc-ch7-2-6}. $\square$
\end{proof}

\begin{theorem}[Hopf and Rinow [HR]]
  \label{thm:dc-ch7-2-8}
  \dcref{ch7:2.8}
  \uses{def:dc-ch7-2-2, def:dc-ch7-2-4, prop:dc-ch7-2-5}
  \lean{Riemannian.Geodesic.hopfRinow}
  \leanok
  Let $M$ be a Riemannian manifold and let $p\in M$. The following assertions are
  equivalent:
  
  a) The intrinsic exponential map $\exp_p$ is defined on all of $T_p(M)$.
  
  b) The closed and bounded sets of $M$ are compact.
  
  c) $M$ is complete as a metric space.
  
  d) $M$ is geodesically complete.
  
  e) There exists a sequence of compact subsets $K_n\subset M$, $K_n\subset K_{n+1}$
     and $\bigcup_n K_n=M$, such that if $q_n\notin K_n$ then $d(p,q_n)\to\infty$.
  
  In addition, any of the statements above implies that
  
  f) For any $q\in M$ there exists a geodesic $\gamma$ joining $p$ to $q$ with
     $\ell(\gamma)=d(p,q)$.
\end{theorem}
\begin{proof}
  \leanok
  \uses{cor:dc-ch3-3-9, lem:dc-ch7-2-8-b-to-c, lem:dc-ch7-2-8-a-to-b, lem:dc-ch7-2-8-b-iff-e, lem:dc-ch7-2-8-lipschitz, lem:dc-ch7-2-8-c-to-d}
  **a) $\Rightarrow$ f).** Let $d(p,q)=r$, and let $B_\delta(p)$ be a normal ball at
  $p$, with $S_\delta(p)=S$ the boundary of $B_\delta(p)$. Let $x_0$ be a point
  where the continuous function $d(q,x)$, $x\in S$, attains a minimum. Then
  $x_0=\exp_p\delta v$, where $v\in T_pM$ and $|v|=1$. Let $\gamma$ be the geodesic
  given by $\gamma(s)=\exp_p sv$. We are going to show that $\gamma(r)=q$.
  
  To prove this fact, consider the equation
  
  $$
  d(\gamma(s),q)=r-s,\qquad\text{(1)}
  $$
  
  and let $A=\{s\in[0,r];\ (1)\text{ is valid}\}$. $A$ is not empty, since (1) is
  true for $s=0$. In addition $A$ is closed in $[0,r]$. Let $s_0\in A$. We show that
  if $s_0<r$, then (1) is valid for $s_0+\delta'$, where $\delta'>0$ is sufficiently
  small. This implies that $\sup A=r$; since $A$ is closed then $r\in A$, which
  shows that $\gamma(r)=q$.
  
  In order to prove that (1) is true for $s_0+\delta'$, let $B_{\delta'}(\gamma(s_0))$
  be a normal ball at $\gamma(s_0)$, with $S'=\partial B_{\delta'}(\gamma(s_0))$ its
  boundary, and let $x_0'$ be a point where $d(x,q)$, $x\in S'$, has a minimum. It
  suffices to show that $x_0'=\gamma(s_0+\delta')$. Indeed, if
  $x_0'=\gamma(s_0+\delta')$, since
  
  $$
  d(\gamma(s_0),q)=\delta'+\min_{x\in S'} d(x,q)=\delta'+d(x_0',q)
  $$
  
  and
  
  $$
  d(\gamma(s_0),q)=r-s_0,
  $$
  
  we have
  
  $$
  r-s_0=\delta'+d(x_0',q)=\delta'+d(\gamma(s_0+\delta'),q),\qquad\text{(2)}
  $$
  
  or that
  
  $$
  d(\gamma(s_0+\delta'),q)=r-(s_0+\delta'),
  $$
  
  which is (1) for $s_0+\delta'$.
  
  To prove finally that $\gamma(s_0+\delta')=x_0'$, observe that, by the triangle
  inequality and by the first equality of (2),
  
  $$
  d(p,x_0')\ge d(p,q)-d(q,x_0')=r-(r-s_0-\delta')=s_0+\delta'.
  $$
  
  On the other hand, the broken curve joining $p$ to $x_0'$ — which goes from $p$ to
  $\gamma(s_0)$ by the geodesic $\gamma$, and from $\gamma(s_0)$ to $x_0'$ by the
  geodesic ray — has length equal to $s_0+\delta'$. Hence $d(p,x_0')=s_0+\delta'$,
  and such a curve, by \cref{cor:dc-ch3-3-9}, is a geodesic. In particular the
  curve is not broken, hence $\gamma(s_0+\delta')=x_0'$. This concludes the proof
  that a) $\Rightarrow$ f).
  
  **a) $\Rightarrow$ b).** Let $A\subset M$ be closed and bounded. Since $A$ is
  bounded, $A\subset B$, where $B$ is a ball with center $p$ in the metric $d$. By
  (f), there exists a ball $B_r(0)\subset T_pM$ such that
  $B\subset\exp_p\overline{B_r(0)}$. Being the continuous image of a compact set,
  $\exp_p\overline{B_r(0)}$ is compact. Hence $A$ is a closed set contained in a
  compact set, and is therefore compact.
  
  **b) $\Rightarrow$ c).** This is \cref{lem:dc-ch7-2-8-b-to-c}.
  
  **c) $\Rightarrow$ d).** Suppose that $M$ is not geodesically complete. Then some
  normalized geodesic $\gamma$ of $M$ is defined for $s<s_0$ and is not defined for
  $s_0$. Let $\{s_n\}$ be a convergent sequence, converging to $s_0$ with $s_n<s_0$.
  Given $\varepsilon>0$, there exists an index $n_0$ such that if $n,m>n_0$ then
  $|s_n-s_m|<\varepsilon$. It follows that
  
  $$
  d(\gamma(s_n),\gamma(s_m))\le|s_n-s_m|<\varepsilon,
  $$
  
  and hence the sequence $\{\gamma(s_n)\}$ is a Cauchy sequence in $M$. Since $M$ is
  complete in the metric $d$, $\{\gamma(s_n)\}\to p_0\in M$.
  
  Let $(W,\delta)$ be a totally normal neighborhood of $p_0$. Choose $n_1$ such that
  if $n,m>n_1$, then $|s_m-s_n|<\delta$ and $\gamma(s_n),\gamma(s_m)$ belong to $W$.
  Then there exists a unique geodesic $g$ whose length is less than $\delta$ joining
  $\gamma(s_n)$ to $\gamma(s_m)$. It is clear that $g$ coincides with $\gamma$,
  wherever $\gamma$ is defined. Since $\exp_{\gamma(s_n)}$ is a diffeomorphism on
  $B_\delta(0)$ and $\exp_{\gamma(s_n)}(B_\delta(0))\supset W$, $g$ extends $\gamma$
  beyond $s_0$. (The formalisation, \cref{lem:dc-ch7-2-8-c-to-d}, replaces the
  totally normal neighborhood of $p_0$ by the uniform-time Picard–Lindelöf flow of
  the geodesic system in the chart at $p_0$, which serves the same purpose: a time
  of existence bounded below uniformly over all initial data near $p_0$.)
  
  **d) $\Rightarrow$ a).** Obvious.

  **b) $\iff$ e).** This is \cref{lem:dc-ch7-2-8-b-iff-e}. $\square$
\end{proof}

\begin{lemma}[proper $\Rightarrow$ complete]
  \label{lem:dc-ch7-2-8-b-to-c}
  \dcref{ch7:2.8}
  \lean{complete_of_proper}
  \mathlibok
  A metric space in which closed bounded sets are compact is complete. (A
  Cauchy sequence is bounded, therefore has compact closure, thus contains a
  convergent subsequence; a Cauchy sequence with a convergent subsequence
  converges.)
\end{lemma}

\begin{lemma}[properness via divergent compact exhaustions]
  \label{lem:dc-ch7-2-8-b-iff-e}
  \dcref{ch7:2.8}
  \lean{OpenGA.properSpace_iff_exists_compact_exhaustion}
  \leanok
  Let $X$ be a metric space and $p\in X$. Closed bounded subsets of $X$ are
  compact if and only if there exists a sequence of compact subsets
  $K_n\subset X$ with $K_n\subset K_{n+1}$ and $\bigcup_n K_n=X$, such that
  $q_n\notin K_n$ implies $d(p,q_n)\to\infty$.
\end{lemma}
\begin{proof}
  \leanok
  If closed bounded sets are compact, take $K_n=\overline{B}_n(p)$, the
  closed metric balls around $p$: they are closed, bounded, increasing, they
  exhaust $X$, and $q_n\notin K_n$ means $d(p,q_n)>n\to\infty$. Conversely,
  given such an exhaustion, every closed ball $\overline{B}_r(p)$ is
  contained in some $K_n$: otherwise for each $n$ we could choose
  $q_n\in\overline{B}_r(p)\smallsetminus K_n$, and then $d(p,q_n)\to\infty$
  while $d(p,q_n)\le r$, a contradiction. Hence $\overline{B}_r(p)$ is a
  closed subset of a compact set, therefore compact. For an arbitrary center
  $x$ one has $\overline{B}_r(x)\subset\overline{B}_{d(p,x)+r}(p)$, so all
  closed balls are compact, and a closed bounded set is a closed subset of
  some closed ball. $\square$
\end{proof}

\begin{lemma}[geodesics are locally Lipschitz]
  \label{lem:dc-ch7-2-8-lipschitz}
  \dcref{ch7:2.8}
  \uses{def:dc-ch3-2-1, lem:dc-ch3-2-1-constspeed, def:dc-ch7-2-4}
  \lean{Riemannian.Geodesic.IsGeodesicOn.contMDiffOn, Riemannian.Geodesic.IsGeodesicOn.edist_le, Riemannian.Geodesic.IsGeodesicOn.dist_le}
  \leanok
  Let $M$ be a Riemannian manifold whose distance $d$ is the Riemannian
  distance of \cref{def:dc-ch7-2-4}, and let $\gamma$ be a (continuous)
  geodesic defined on a connected open set $I$ of times. Then $\gamma$ is
  continuously differentiable, and for $a\le b$ in $I$,
  $$
  d(\gamma(a),\gamma(b))\ \le\ |\gamma'|\,(b-a),
  $$
  where $|\gamma'|=\sqrt{\langle\gamma',\gamma'\rangle}$ is the (constant)
  speed of $\gamma$ (\cref{lem:dc-ch3-2-1-constspeed}). In particular a
  normalized geodesic satisfies $d(\gamma(s),\gamma(t))\le|s-t|$: geodesics
  are locally Lipschitz. This is the inequality behind the Cauchy-sequence
  step of the proof of \cref{thm:dc-ch7-2-8}, c) $\Rightarrow$ d).
\end{lemma}
\begin{proof}
  \leanok
  \uses{lem:dc-ch3-2-1-constspeed}
  By \cref{lem:dc-ch3-2-1-constspeed} the squared speed
  $\langle\gamma'(t),\gamma'(t)\rangle$ is constant on $I$, say equal to
  $c^2$ with $c\ge 0$. The chart velocity of $\gamma$ transforms under a
  change of chart by the derivative of the smooth transition map, which
  depends continuously on the point; combined with the differentiability
  supplied by the geodesic equation this makes $\gamma$ of class $C^1$ on
  $I$. The length of $\gamma$ over $[a,b]\subset I$ is therefore defined and
  equals $\int_a^b|\gamma'(t)|\,dt=c\,(b-a)$. Since $d$ is the infimum of the
  lengths of piecewise-differentiable curves joining the endpoints
  (\cref{def:dc-ch7-2-4}), the restriction of $\gamma$ to $[a,b]$ witnesses
  $d(\gamma(a),\gamma(b))\le c\,(b-a)$. $\square$
\end{proof}

\begin{lemma}[transformation law for the Christoffel symbols]
  \label{lem:dc-ch7-2-8-christoffel-transform}
  \dcref{ch7:2.8}
  \uses{def:dc-ch2-3-1, thm:dc-ch2-3-6}
  \lean{Riemannian.chartChristoffelContraction_change}
  \leanok
  Let $\varphi_\alpha,\varphi_\beta$ be two charts of $M$ whose domains contain a
  point $x$, let $\tau=\varphi_\beta\circ\varphi_\alpha^{-1}$ be the transition map,
  $A=D\tau(\varphi_\alpha(x))$ its derivative and $D^2\tau$ its second derivative at
  $\varphi_\alpha(x)$. Then the Christoffel contractions
  $\Gamma^\alpha,\Gamma^\beta$ of the Levi–Civita connection of $g$ in the two
  charts satisfy, for all $v,w\in\mathbb R^n$,
  $$
  A\,\Gamma^\alpha(v,w)\;=\;\Gamma^\beta(Av,Aw)\;+\;D^2\tau(v,w).
  $$
\end{lemma}
\begin{proof}
  \leanok
  Write $y$ for $\alpha$-coordinates, $x=\tau(y)$ for $\beta$-coordinates,
  $A^a_i(y)=\partial\tau^a/\partial y^i$ and
  $B^a_{ki}=\partial^2\tau^a/\partial y^k\partial y^i$ (symmetric in $k,i$ by
  Schwarz). The Gram matrices of $g$ in the two charts are related by
  $G^\alpha_{ij}=A^a_iA^b_j\,(G^\beta_{ab}\circ\tau)$. Differentiating in the
  direction $k$,
  $$
  \partial_kG^\alpha_{ij}
   = B^a_{ki}A^b_jG^\beta_{ab}+A^a_iB^b_{kj}G^\beta_{ab}
     +A^a_iA^b_jA^c_k\,\partial_cG^\beta_{ab}.
  $$
  Form the first-kind bracket
  $[ij,k]:=\tfrac12(\partial_iG_{jk}+\partial_jG_{ik}-\partial_kG_{ij})$. Inserting
  the displayed identity three times and using the symmetry of $B$ in its lower
  indices and of $G^\beta$, the six $B$-terms cancel in pairs except for
  $G^\beta_{ab}B^a_{ij}A^b_k$, and the three $\partial G^\beta$-terms combine to the
  $\beta$-bracket:
  $$
  [ij,k]^\alpha=G^\beta_{ab}\,B^a_{ij}\,A^b_k+A^a_iA^b_jA^c_k\,[ab,c]^\beta .
  $$
  The Christoffel symbols are characterised by the contraction identity
  $\sum_mG_{mk}\Gamma^m_{ij}=[ij,k]$ (\cref{thm:dc-ch2-3-6}). Substituting it on
  both sides, and writing $G^\alpha_{mk}=A^c_mA^b_k(G^\beta_{cb}\circ\tau)$ on the
  left,
  $$
  G^\beta_{cb}\,\bigl(A^c_m\Gamma^{\alpha,m}_{ij}\bigr)A^b_k
   = G^\beta_{cb}\,\bigl(B^c_{ij}
     +\Gamma^{\beta,c}_{ab}A^a_iA^b_j\bigr)A^b_k .
  $$
  The factor $A^b_k$ ranges over the columns of the invertible matrix $A$ and
  $G^\beta$ is invertible (positive definite), so the parenthesised vectors agree:
  $A^c_m\Gamma^{\alpha,m}_{ij}=B^c_{ij}+\Gamma^{\beta,c}_{ab}A^a_iA^b_j$.
  Contracting with $v^iw^j$ gives the claim. $\square$
\end{proof}

\begin{lemma}[the geodesic equation is chart-independent]
  \label{lem:dc-ch7-2-8-chart-independence}
  \dcref{ch7:2.8}
  \uses{def:dc-ch3-2-1, lem:dc-ch7-2-8-christoffel-transform}
  \lean{Riemannian.Geodesic.SolvesGeodesicODEAt, Riemannian.Geodesic.SolvesGeodesicODEAt.transfer, Riemannian.Geodesic.HasGeodesicEquationAt.solvesGeodesicODEAt, Riemannian.Geodesic.SolvesGeodesicODEAt.hasGeodesicEquationAt}
  \leanok
  Let $\gamma$ be a curve in $M$ continuous at $t_0$ whose foot $\gamma(t_0)$ lies
  in the domains of two charts $\varphi_\alpha,\varphi_\beta$, and write
  $u_\alpha=\varphi_\alpha\circ\gamma$, $u_\beta=\varphi_\beta\circ\gamma$. If
  $u_\alpha$ is differentiable near $t_0$, its derivative is differentiable at
  $t_0$, and the second-order geodesic equation
  $u_\alpha''+\Gamma^\alpha(u_\alpha',u_\alpha')(u_\alpha)=0$ holds at $t_0$, then
  the same three statements hold for $u_\beta$. In particular the geodesic
  equation of \cref{def:dc-ch3-2-1}, read in the chart at the moving foot
  $\gamma(t_0)$, holds if and only if it holds in any fixed chart whose domain
  contains the foot.
\end{lemma}
\begin{proof}
  \leanok
  \uses{lem:dc-ch7-2-8-christoffel-transform}
  By continuity of $\gamma$ at $t_0$ the foot stays in both chart domains for
  times near $t_0$, so $u_\beta=\tau\circ u_\alpha$ near $t_0$ with
  $\tau=\varphi_\beta\circ\varphi_\alpha^{-1}$ smooth near $u_\alpha(t_0)$. The
  chain rule gives $u_\beta'=D\tau(u_\alpha)\,u_\alpha'$ near $t_0$ — hence
  differentiability of $u_\beta$ near $t_0$ — and, differentiating once more at
  $t_0$ with $v=u_\alpha'(t_0)$, $A=D\tau(u_\alpha(t_0))$,
  $$
  u_\beta''(t_0)=D^2\tau(v,v)+A\,u_\alpha''(t_0).
  $$
  Adding $\Gamma^\beta(u_\beta',u_\beta')(u_\beta)=\Gamma^\beta(Av,Av)$ at $t_0$ and
  using the transformation law of \cref{lem:dc-ch7-2-8-christoffel-transform},
  $$
  u_\beta''+\Gamma^\beta(u_\beta',u_\beta')
   = D^2\tau(v,v)+A u_\alpha''+A\,\Gamma^\alpha(v,v)-D^2\tau(v,v)
   = A\bigl(u_\alpha''+\Gamma^\alpha(v,v)\bigr)=0 .
  $$
  Since $A$ is invertible, the situation is symmetric in $\alpha,\beta$; taking
  $\varphi_\beta$ to be the canonical chart at the foot identifies the fixed-chart
  equation with the moving-foot equation of \cref{def:dc-ch3-2-1}. $\square$
\end{proof}

\begin{lemma}[flow solutions are geodesics; uniform local existence]
  \label{lem:dc-ch7-2-8-flow-geodesic}
  \dcref{ch7:2.8}
  \uses{def:dc-ch3-2-1, lem:dc-ch7-2-8-chart-independence, lem:dc-ch3-2-2-pointwise, lem:dc-ch3-2-6-reparam}
  \lean{Riemannian.Geodesic.sprayBase, Riemannian.Geodesic.isGeodesicOn_sprayBase, Riemannian.Geodesic.hasDerivAt_chartReading_sprayBase, Riemannian.Geodesic.continuousOn_sprayBase, Riemannian.Geodesic.exists_seed_geodesic}
  \leanok
  Let $q\in M$ and let $\zeta=(x,w):J\to\mathbb R^n\times\mathbb R^n$ solve the
  first-order geodesic system $x'=w$, $w'=-\Gamma^q(w,w)(x)$ of the chart at $q$
  on an open set $J$ of times, with $x$ staying in the chart image. Then the base
  curve $\sigma=\varphi_q^{-1}\circ x$ is continuous and satisfies the (intrinsic)
  geodesic equation at every time of $J$, with chart velocity reading $w$.
  Consequently — by Picard–Lindelöf applied to this system, whose right-hand side
  is $C^1$ on the chart image — every initial datum $(p,v)$ admits, for some
  $b>0$, a continuous geodesic on $(-b,b)$ through $p$ with chart-$p$ velocity
  $v$ at time $0$.
\end{lemma}
\begin{proof}
  \leanok
  \uses{lem:dc-ch7-2-8-chart-independence, lem:dc-ch3-2-6-reparam}
  $x$ is differentiable with $x'=w$ and $w$ is differentiable with
  $w'=-\Gamma^q(w,w)(x)$, so $u:=x=\varphi_q\circ\sigma$ satisfies
  $u''+\Gamma^q(u',u')(u)=0$ at every time of $J$; $\sigma$ is continuous as
  $\varphi_q^{-1}$ is. By \cref{lem:dc-ch7-2-8-chart-independence} the fixed-chart
  equation upgrades to the moving-foot geodesic equation at every $t\in J$.
  For the last claim, apply Picard–Lindelöf to the system at the initial condition
  $(\varphi_p(p),\kappa_0v)$, where $\kappa_0>0$ is chosen with
  $\|\kappa_0v\|\le r$ for the Picard–Lindelöf radius $r$ of the chart at $p$;
  this yields a geodesic $\sigma$ on some $(-\varepsilon,\varepsilon)$ through $p$
  with velocity $\kappa_0v$, and the affine rescaling
  $t\mapsto\sigma(\kappa_0^{-1}t)$ (\cref{lem:dc-ch3-2-6-reparam}) is a geodesic
  on $(-\kappa_0\varepsilon,\kappa_0\varepsilon)$ through $p$ with velocity $v$.
  $\square$
\end{proof}

\begin{lemma}[uniqueness of geodesics with given initial data]
  \label{lem:dc-ch7-2-8-uniqueness}
  \dcref{ch7:2.8}
  \uses{def:dc-ch3-2-1, lem:dc-ch7-2-8-chart-independence}
  \lean{Riemannian.Geodesic.IsGeodesicOn.eventuallyEq_of_deriv_chartReading_eq, Riemannian.Geodesic.IsGeodesicOn.eqOn_of_deriv_chartReading_eq}
  \leanok
  Let $M$ be Hausdorff and let $\gamma_1,\gamma_2$ be continuous geodesics on an
  open interval $I$ which at some $t_0\in I$ have the same value and the same
  velocity reading in some chart around the common foot. Then
  $\gamma_1=\gamma_2$ on $I$. (Hausdorffness is necessary: on the line with two
  origins, two geodesics can agree except at a single time, where they pass
  through the two origins.)
\end{lemma}
\begin{proof}
  \leanok
  \uses{lem:dc-ch7-2-8-chart-independence}
  \emph{Local step.} Near a time where the curves share value and chart velocity,
  read both in a fixed chart at the common foot. By
  \cref{lem:dc-ch7-2-8-chart-independence} the two chart readings $u_1,u_2$
  solve the second-order geodesic equation near that time, so their first-order
  lifts $(u_i,u_i')$ solve the system $x'=w$, $w'=-\Gamma(w,w)(x)$, whose
  right-hand side is $C^1$, hence locally Lipschitz. Grönwall's inequality forces
  the lifts to agree on a neighbourhood, and injectivity of the chart returns
  $\gamma_1=\gamma_2$ there.

  \emph{Propagation.} Let $S\subset I$ be the set of times near which the curves
  eventually agree; $S$ is open and contains $t_0$. If $t^*\in I$ is adherent to
  $S$, continuity and Hausdorffness give $\gamma_1(t^*)=\gamma_2(t^*)$ (limits
  along the approach times are unique), and the difference of the chart readings
  at the common foot is differentiable at $t^*$ and vanishes on a set of times
  accumulating at $t^*$, so its derivative — a limit of vanishing difference
  quotients — is zero: the velocities also match at $t^*$. The local step then
  puts $t^*\in S$, so $S$ is closed in $I$; by connectedness $S=I$. $\square$
\end{proof}

\begin{lemma}[local comparison of coordinate and Riemannian norms]
  \label{lem:dc-ch7-2-8-gram-bound}
  \dcref{ch7:2.8}
  \uses{def:dc-ch1-2-1, lem:dc-ch3-2-1-constspeed}
  \lean{Riemannian.Geodesic.exists_sq_norm_le_chartMetricInner, Riemannian.Geodesic.HasGeodesicEquationAt.speedSq_eq_chartMetricInner_of_mem_source}
  \leanok
  Let $p_0\in M$ and let $G$ be the Gram matrix of $g$ in the chart at $p_0$.
  There are $c>0$ and a neighbourhood $V$ of $\varphi_{p_0}(p_0)$ inside the chart
  image such that $\|w\|^2\le c\,\langle w,w\rangle_{G(y)}$ for all $y\in V$ and
  all $w\in\mathbb R^n$. Moreover, along any geodesic whose foot lies in the chart
  at $p_0$, the Gram reading $\langle u',u'\rangle_{G(u)}$ of the chart velocity
  equals the intrinsic squared speed $\langle\gamma',\gamma'\rangle$.
\end{lemma}
\begin{proof}
  \leanok
  At $y_0=\varphi_{p_0}(p_0)$ the Gram form is the inner product $g_{p_0}$ itself,
  hence positive definite; so on the compact unit sphere the continuous function
  $w\mapsto\langle w,w\rangle_{G(y_0)}$ attains a minimum $m>0$. The map
  $(y,w)\mapsto\langle w,w\rangle_{G(y)}$ is continuous, so
  $\{(y,w):\langle w,w\rangle_{G(y)}>m/2\}$ is open and contains
  $\{y_0\}\times S^{n-1}$; by the tube lemma it contains $V\times S^{n-1}$ for
  some neighbourhood $V\ni y_0$. Homogeneity of degree $2$ in $w$ then gives
  $\langle w,w\rangle_{G(y)}\ge\tfrac m2\|w\|^2$ on $V$, i.e. the claim with
  $c=2/m$. For the second statement, the chart velocity reading of a geodesic at
  a time $\sigma$ is the coordinate change of the moving-foot velocity, and the
  Gram matrix at the foot represents $g$ there, so
  $\langle u'(\sigma),u'(\sigma)\rangle_{G(u(\sigma))}
   =g_{\gamma(\sigma)}(\gamma'(\sigma),\gamma'(\sigma))$. $\square$
\end{proof}

\begin{lemma}[c $\Rightarrow$ d: complete metrics have global geodesics]
  \label{lem:dc-ch7-2-8-c-to-d}
  \dcref{ch7:2.8}
  \uses{def:dc-ch7-2-2, def:dc-ch7-2-4, lem:dc-ch7-2-8-lipschitz, lem:dc-ch7-2-8-flow-geodesic, lem:dc-ch7-2-8-uniqueness, lem:dc-ch7-2-8-gram-bound, lem:dc-ch3-2-6-reparam}
  \lean{Riemannian.Geodesic.exists_tendsto_of_isGeodesicOn, Riemannian.Geodesic.IsGeodesicOn.exists_forward_extension, Riemannian.Geodesic.exists_global_geodesic, Riemannian.Geodesic.isGeodesicallyComplete_of_complete}
  \leanok
  Let $M$ be a Riemannian manifold whose distance is the Riemannian distance of
  $g$ and which is complete as a metric space. Then for every $p\in M$ and
  $v\in T_pM$ there is a geodesic $\gamma:\mathbb R\to M$ with $\gamma(0)=p$ and
  $\gamma'(0)=v$: $M$ is geodesically complete.
\end{lemma}
\begin{proof}
  \leanok
  \uses{lem:dc-ch7-2-8-lipschitz, lem:dc-ch7-2-8-flow-geodesic, lem:dc-ch7-2-8-uniqueness, lem:dc-ch7-2-8-gram-bound, lem:dc-ch3-2-6-reparam}
  \emph{Extension step.} Let $\gamma$ be a continuous geodesic on $(a,b)$ with
  $b<\infty$. Its speed is a constant $\sqrt S$, so by
  \cref{lem:dc-ch7-2-8-lipschitz} $d(\gamma(s),\gamma(t))\le\sqrt S\,|t-s|$: the
  values $\gamma(\tau)$ form a Cauchy filter as $\tau\to b^-$ and converge to some
  $p_0\in M$ by completeness. Fix the Picard–Lindelöf data of the chart at $p_0$:
  a radius $r>0$ and a uniform time $\varepsilon>0$ such that the geodesic system
  of \cref{lem:dc-ch7-2-8-flow-geodesic} is solvable on $(-\varepsilon,\varepsilon)$
  from every initial condition $r$-close to $(\varphi_{p_0}(p_0),0)$, and the Gram
  comparison constant $c$ of \cref{lem:dc-ch7-2-8-gram-bound}. Near $b$ the foot
  of $\gamma$ lies in the chart at $p_0$, and its chart velocity is bounded by
  $\sqrt{cS}$ by \cref{lem:dc-ch7-2-8-gram-bound} and constancy of speed. Choose
  $\kappa\in(0,1]$ with $\kappa\sqrt{cS}\le r$; the affine reparametrisation
  $\delta(t)=\gamma(b+\kappa(t-b))$ (\cref{lem:dc-ch3-2-6-reparam}) is a geodesic
  with the same endpoint behaviour and chart velocity bounded by $r$. Pick
  $t_1<b$ with $b-t_1<\varepsilon/2$, with $\delta(t_1)$ in the chart at $p_0$ and
  its chart position $r$-close to $\varphi_{p_0}(p_0)$. The flow solution through
  the chart data of $\delta$ at $t_1$ is a geodesic on
  $(t_1-\varepsilon,t_1+\varepsilon)$, which by
  \cref{lem:dc-ch7-2-8-uniqueness} agrees with $\delta$ on their common interval;
  since $t_1+\varepsilon>b+\varepsilon/2$, the glued curve prolongs $\delta$ —
  and, undoing the reparametrisation, prolongs $\gamma$ to
  $(a,b+\kappa\varepsilon/2)$.

  \emph{Global geodesic.} Given $(p,v)$, consider the set $\mathcal S$ of radii
  $b>0$ such that some continuous geodesic on $(-b,b)$ has data $(p,v)$ at $0$;
  it is nonempty by \cref{lem:dc-ch7-2-8-flow-geodesic}. Any two members agree on
  the intersection of their intervals by \cref{lem:dc-ch7-2-8-uniqueness}, so the
  family glues to a single geodesic on $(-T,T)$, $T=\sup\mathcal S$. If
  $T<\infty$, the extension step — applied forward, and applied to the
  time-reversed curve for the backward direction — produces a geodesic on a
  strictly larger symmetric interval with the same data, contradicting the
  definition of $T$. Hence $T=\infty$ and the glued curve is the required global
  geodesic. $\square$
\end{proof}

\begin{lemma}[a $\Rightarrow$ b: geodesic completeness at a point makes closed balls compact]
  \label{lem:dc-ch7-2-8-a-to-b}
  \dcref{ch7:2.8}
  \uses{def:dc-ch7-2-2, def:dc-ch7-2-4, lem:dc-ch7-2-8-step}
  \lean{Riemannian.Geodesic.properSpace_of_expDomainIntrinsic_eq_univ}
  \leanok
  Let $M$ be a connected Riemannian manifold whose distance is the Riemannian distance of
  $g$, and let $p\in M$ be a point at which $M$ is geodesically complete: every $v\in T_pM$
  generates a geodesic $\gamma:\mathbb R\to M$, defined for all time, with $\gamma(0)=p$ and
  $\gamma'(0)=v$. Then the closed and bounded subsets of $M$ are compact, i.e. $M$ is a
  \emph{proper} metric space. This is do Carmo's a) $\Rightarrow$ b) branch of the
  Hopf--Rinow theorem read at a single base point; together with
  \cref{lem:dc-ch7-2-8-b-to-c} it gives the metric-completeness form
  (\cref{thm:dc-ch7-2-8}, a) $\Rightarrow$ c)). It is the properness invoked in the proof of
  the Hadamard theorem for $T_pM$ equipped with the pulled-back metric.
\end{lemma}
\begin{proof}
  \leanok
  \uses{def:dc-ch7-2-4, lem:dc-ch7-2-8-step}
  Let $A\subset M$ be closed and bounded; then $A\subset\overline B_R(p)$ for some
  $R<\infty$. By the a) $\Rightarrow$ f) growth argument (one step of which is
  \cref{lem:dc-ch7-2-8-step}) every point of the closed ball $\overline B_R(p)$ is joined to
  $p$ by a minimizing geodesic, hence is $\exp_p(w)$ for some $w\in T_pM$ with $|w|_p\le R$;
  so $\overline B_R(p)\subset\exp_p\big(\overline{B_R(0)}\big)$. Being the continuous image
  of the compact Euclidean ball $\overline{B_R(0)}\subset T_pM$, the set
  $\exp_p\big(\overline{B_R(0)}\big)$ is compact. Hence $A$ is a closed subset of a compact
  set, and is therefore compact. $\square$
\end{proof}

\begin{lemma}[length of a curve read in one chart]
  \label{lem:dc-ch7-2-8-chart-length}
  \dcref{ch7:2.8}
  \uses{def:dc-ch7-2-4, lem:dc-ch7-2-8-gram-bound}
  \lean{Riemannian.Geodesic.hasMFDerivAt_of_hasDerivAt_extChartAt, Riemannian.Geodesic.mfderiv_eq_of_hasDerivAt_extChartAt, Riemannian.Geodesic.enorm_mfderiv_eq_of_hasDerivAt_extChartAt, Riemannian.Geodesic.contDiffOn_extChartAt_comp, Riemannian.Geodesic.pathELength_eq_ofReal_integral_chartMetricInner}
  \leanok
  Let $\gamma:[a,b]\to M$ be a $C^1$ curve whose image lies in the domain of a
  single chart $\varphi_\alpha$, and let $u=\varphi_\alpha\circ\gamma$ be its
  chart reading. Then at every $t$ the intrinsic velocity $\gamma'(t)\in
  T_{\gamma(t)}M$ is the readback of the coordinate velocity $\dot u(t)$ under
  the chart frame at $\gamma(t)$, its Riemannian norm is the Gram norm
  $|\gamma'(t)| = \sqrt{\langle\dot u(t),\dot u(t)\rangle_{G(u(t))}}$ of the
  coordinate velocity, and the length of $\gamma$ (\cref{def:dc-ch7-2-4}) is
  the coordinate integral
  $$
  \ell(\gamma)\;=\;\int_a^b \sqrt{\big\langle\dot u(t),\dot u(t)\big\rangle_{G(u(t))}}\,dt ,
  $$
  where $G$ is the Gram matrix of $g$ in the chart at $\alpha$.
\end{lemma}
\begin{proof}
  \leanok
  \uses{lem:dc-ch7-2-8-gram-bound}
  Fix $t$ and read the tangent space $T_{\gamma(t)}M$ in the chart at
  $\gamma(t)$ itself. Near $t$ the reading of $\gamma$ in the chart at
  $\gamma(t)$ is the composition of the smooth transition
  $\varphi_{\gamma(t)}\circ\varphi_\alpha^{-1}$ with $u$; by the chain rule its
  derivative at $t$ is $A\,\dot u(t)$, where
  $A=d(\varphi_{\gamma(t)}\circ\varphi_\alpha^{-1})_{u(t)}$ is the tangent
  coordinate change. This says exactly that $\gamma$ is differentiable at $t$
  as a curve into $M$ with velocity the readback of $\dot u(t)$ through the
  $\alpha$-chart frame of $T_{\gamma(t)}M$. Since the Gram matrix at the foot
  represents $g$ there (\cref{lem:dc-ch7-2-8-gram-bound}), the Riemannian norm
  of this readback is
  $\sqrt{\langle\dot u(t),\dot u(t)\rangle_{G(u(t))}}$. The chart reading $u$
  is $C^1$ on $[a,b]$ (composition of the $C^1$ curve with the smooth chart),
  so this integrand is continuous, and integrating the pointwise identity over
  $[a,b]$ — the length of \cref{def:dc-ch7-2-4} being the integral of
  $|\gamma'|$ — gives the displayed formula. $\square$
\end{proof}

\begin{lemma}[the metric normal ball]
  \label{lem:dc-ch7-2-8-normal-ball-dist}
  \dcref{ch7:2.8}
  \uses{def:dc-ch7-2-4, def:dc-ch3-expmap, lem:dc-ch3-2-9-c1diffeo, lem:dc-ch3-3-6-radius, lem:dc-ch3-3-6-reach, lem:dc-ch3-3-6-rayspeed, lem:dc-ch7-2-8-gram-bound, lem:dc-ch7-2-8-chart-length}
  \lean{Riemannian.Exponential.exists_isOpen_expMap_image, Riemannian.Exponential.exists_pathELength_expMap_ray, Riemannian.Exponential.exists_le_pathELength, Riemannian.Exponential.exists_edist_expMap_ball}
  \leanok
  Let $M$ be a Riemannian manifold whose distance $d$ is the Riemannian
  distance of \cref{def:dc-ch7-2-4} and let $p\in M$. There are
  $\varepsilon,\delta>0$ such that $\exp_p$ is injective and open on
  $B_\varepsilon(0)\subset T_pM$ and:
  \begin{enumerate}
  \item (\emph{radial geodesics realize the distance})
    $d\big(p,\exp_p v\big)=|v|_p$ for every $v$ with $\|v\|<\varepsilon$;
  \item (\emph{escape estimate}) $d(p,q)\ge\delta$ for every
    $q\notin\exp_p\big(B_\varepsilon(0)\big)$ — the normal ball contains the
    metric ball $B_\delta(p)$.
  \end{enumerate}
  This is the metric form of \cref{prop:dc-ch3-3-6} consumed by the proof of
  \cref{thm:dc-ch7-2-8}: inside a normal ball, distances to the center are
  computed by radial geodesics, and leaving the normal ball costs a definite
  length.
\end{lemma}
\begin{proof}
  \leanok
  \uses{lem:dc-ch3-2-9-c1diffeo, lem:dc-ch3-3-6-radius, lem:dc-ch3-3-6-reach, lem:dc-ch3-3-6-rayspeed, lem:dc-ch7-2-8-gram-bound, lem:dc-ch7-2-8-chart-length}
  Choose $\rho>0$ small enough that $B_\rho(0)$ lies in the domains of
  the $C^1$ diffeomorphism property of $\exp_p$
  (\cref{lem:dc-ch3-2-9-c1diffeo}, with local inverse $\exp_p^{-1}$ on the
  image), of the radius comparison (\cref{lem:dc-ch3-3-6-radius}), and of the
  radial speed identity (\cref{lem:dc-ch3-3-6-rayspeed}); let $c>0$ be the
  Gram comparison constant of \cref{lem:dc-ch7-2-8-gram-bound} at $p$, so that
  $\|w\|/\sqrt c\le|w|_p$ for all $w\in T_pM$. Set $r'=\rho/2$ and
  $U=\exp_p(B_{r'}(0))$, an open set containing $p$. Openness of the image of
  any open $V\subset B_\rho(0)$ follows because the chart reading of $\exp_p$
  has an invertible derivative at every point of $B_\rho(0)$
  (\cref{lem:dc-ch3-2-9-c1diffeo}), hence is an open map, and the chart is a
  homeomorphism onto its image.

  \emph{Upper bound.} For $\|v\|<\rho$ the radial curve
  $\gamma_v(t)=\exp_p(tv)$, $t\in[0,1]$, is $C^1$ (its chart reading is the
  $C^1$ map $w\mapsto\varphi_p(\exp_p w)$ composed with $t\mapsto tv$), stays
  in the chart at $p$, and by \cref{lem:dc-ch3-3-6-rayspeed} its chart-read
  speed is constantly $|v|_p$; by \cref{lem:dc-ch7-2-8-chart-length} its
  length is $|v|_p$. Hence $d(p,\exp_p v)\le|v|_p$.

  \emph{Competitor bound.} Let $\sigma:[0,1]\to M$ be any $C^1$ curve with
  $\sigma(0)=p$. Two cases.

  (i) \emph{$\sigma$ stays in $U$.} Then the polar lift
  $w=\exp_p^{-1}\circ\sigma$ is a $C^1$ curve in
  $B_{r'}(0)$ with $w(0)=0$ and $\exp_p(w(t))=\sigma(t)$; differentiating this
  identity, the chart velocity of $\sigma$ is
  $d(\varphi_p\circ\exp_p)_{w(t)}\dot w(t)$.
  The radius comparison (\cref{lem:dc-ch3-3-6-radius}) applied to $w$ on
  $[0,1]$, combined with \cref{lem:dc-ch7-2-8-chart-length}, gives
  $$
  |w(1)|_p \;=\; |w(1)|_p-|w(0)|_p\;\le\;\ell(\sigma).
  $$

  (ii) \emph{$\sigma$ leaves $U$.} Let $T=\inf\{t:\sigma(t)\notin U\}$; then
  $T>0$, $\sigma(t)\in U$ for $t<T$, and $\sigma(T)$ lies in the compact set
  $\exp_p(\overline B_{r'}(0))$ (the continuous image of a closed ball, which
  is closed and contains all $\sigma(t)$, $t<T$) but not in the open set $U$;
  hence $\sigma(T)=\exp_p(z_0)$ with $\|z_0\|=r'$ exactly. The reach estimate
  (\cref{lem:dc-ch3-3-6-reach}) applied to the polar lift on $[0,T]$, with
  \cref{lem:dc-ch7-2-8-chart-length} and \cref{lem:dc-ch7-2-8-gram-bound},
  gives
  $$
  \ell(\sigma)\;\ge\;\ell\big(\sigma|_{[0,T]}\big)\;\ge\;|z_0|_p
  \;\ge\;\frac{\|z_0\|}{\sqrt c}\;=\;\frac{r'}{\sqrt c}.
  $$

  \emph{Conclusion.} By continuity of the Gram form at $0$ choose
  $\varepsilon\le r'$ so small that $\|v\|<\varepsilon$ implies
  $|v|_p<r'/\sqrt c$, and set $\delta=\varepsilon/\sqrt c$. For
  $\|v\|<\varepsilon$, every competitor from $p$ to $\exp_p v$ is, by (i) with
  $w(1)=v$ (injectivity of $\exp_p$) or by (ii) together with
  $|v|_p<r'/\sqrt c$, at least $|v|_p$ long; taking the infimum,
  $d(p,\exp_p v)=|v|_p$. For $q\notin\exp_p(B_\varepsilon(0))$, every curve
  from $p$ to $q$ either leaves $U$ — costing $r'/\sqrt c\ge\delta$ by (ii) —
  or stays in $U$ and ends at $\exp_p(w(1))$ with $\|w(1)\|\ge\varepsilon$,
  costing $|w(1)|_p\ge\varepsilon/\sqrt c=\delta$ by (i); hence
  $d(p,q)\ge\delta$. $\square$
\end{proof}

\begin{lemma}[exponential rays are geodesics]
  \label{lem:dc-ch7-2-8-ray-geodesic}
  \dcref{ch7:2.8}
  \uses{def:dc-ch3-expmap, def:dc-ch3-2-1, lem:dc-ch7-2-8-chart-independence}
  \lean{Riemannian.Exponential.exists_isGeodesicOn_expMap_ray}
  \leanok
  Let $p\in M$. There are $\rho>0$ and $b>1$ such that for every $u\in T_pM$
  with $\|u\|<\rho$ the ray
  $$
  \gamma_u : t \longmapsto \exp_p(t\,u)
  $$
  is defined and stays in the chart at $p$ for $|t|<b$, starts at
  $\gamma_u(0)=p$ with velocity $u$, is continuous, and satisfies the geodesic
  equation (\cref{def:dc-ch3-2-1}) at every $t\in(-b,b)$. In particular the
  segment $t\in[0,1]$, joining $p$ to $\exp_p u$, is a geodesic segment.
\end{lemma}
\begin{proof}
  \leanok
  \uses{lem:dc-ch7-2-8-chart-independence}
  Read everything in the chart at $p$ and write $f=\varphi_p\circ\exp_p$ for
  the chart reading of the exponential map. On a small ball $B_\rho(0)$ the
  map $f$ is $C^2$, $df_0=\mathrm{id}$, and along rays the velocity
  $t\mapsto df_{tu}(u)$ satisfies the second-order chart geodesic equation
  $$
  \frac{d}{dt}\,df_{tu}(u)\;=\;-\,\Gamma_p\big(df_{tu}(u),\,df_{tu}(u)\big)(f(tu)),
  $$
  because $\exp_p(tu)$ is by construction the time-$t$ point of the geodesic
  flow with initial datum $(p,u)$, suitably rescaled (homogeneity). By the
  chain rule the chart reading $t\mapsto f(tu)$ of $\gamma_u$ is $C^1$ with
  derivative $df_{tu}(u)$, so the displayed identity says exactly that this
  reading solves the fixed-chart geodesic ODE at every $|t|<b$; its value at
  $t=0$ is $p$ and its derivative there is $df_0(u)=u$. Continuity of
  $\gamma_u$ follows by composing with the inverse chart. Finally, by the
  chart-independence of the geodesic equation
  (\cref{lem:dc-ch7-2-8-chart-independence}), a continuous curve solving the
  fixed-chart ODE in a chart containing its points satisfies the intrinsic
  geodesic equation of \cref{def:dc-ch3-2-1}. $\square$
\end{proof}

\begin{lemma}[sphere-minimum decomposition of the distance]
  \label{lem:dc-ch7-2-8-sphere-min}
  \dcref{ch7:2.8}
  \uses{def:dc-ch7-2-4, lem:dc-ch7-2-8-gram-bound, lem:dc-ch7-2-8-chart-length, lem:dc-ch7-2-8-normal-ball-dist}
  \lean{Riemannian.Exponential.exists_normalSphere_min_edist}
  \leanok
  Let $M$ be a Riemannian manifold whose distance $d$ is the Riemannian
  distance of \cref{def:dc-ch7-2-4} and let $p\in M$. There are
  $\varepsilon, c>0$, depending only on $p$, with the following property: for
  every $\delta>0$ with $\sqrt c\,\delta<\varepsilon$, writing
  $S_\delta(p)=\exp_p\{\,|z|_p=\delta\,\}$ for the geodesic sphere of radius
  $\delta$, and for every $q\in M$ with $d(p,q)\ge\delta$, there is a point
  $x_0=\exp_p z_0\in S_\delta(p)$ such that
  $$
  d(p,q)\;=\;\delta+d(x_0,q),
  \qquad
  d(x_0,q)\;=\;\min_{x\in S_\delta(p)} d(x,q).
  $$
  This is the step of do Carmo's proof of \cref{thm:dc-ch7-2-8} that reads
  ``since $d(p,q)=\delta+\min_{x\in S_\delta}d(x,q)$'': every path from $p$ to
  a point beyond the geodesic ball first crosses the geodesic sphere, and the
  crossing costs exactly $\delta$.
\end{lemma}
\begin{proof}
  \leanok
  \uses{lem:dc-ch7-2-8-gram-bound, lem:dc-ch7-2-8-chart-length, lem:dc-ch7-2-8-normal-ball-dist}
  Take $\varepsilon$ small enough that the conclusions of
  \cref{lem:dc-ch7-2-8-normal-ball-dist} hold on $B_\varepsilon(0)\subset
  T_pM$ and let $c$ be the Gram comparison constant of
  \cref{lem:dc-ch7-2-8-gram-bound} at $p$, so $\|z\|\le\sqrt c\,|z|_p$ for all
  $z$. Fix $\delta$ with $0<\sqrt c\,\delta<\varepsilon$.

  \emph{The sphere is compact and nonempty.} The set
  $\Sigma=\{z\in T_pM: |z|_p=\delta\}\cap\overline B_{\sqrt c\,\delta}(0)$ is
  closed and bounded, hence compact, and by the Gram bound every $z$ with
  $|z|_p=\delta$ has $\|z\|\le\sqrt c\,\delta<\varepsilon$, so
  $S_\delta(p)=\exp_p(\Sigma)$; it is nonempty because $z\mapsto|z|_p$ is
  continuous, positive on any unit vector $w_0$ (again by the Gram bound), and
  scales linearly, so $z_\delta=(\delta/|w_0|_p)\,w_0\in\Sigma$. Since
  $\exp_p$ is continuous on $\overline B_{\sqrt c\,\delta}(0)$ and
  $x\mapsto d(x,q)$ is continuous, the minimum
  $m=\min_{x\in S_\delta(p)}d(x,q)$ is attained, say at $x_0=\exp_p z_0$,
  $|z_0|_p=\delta$.

  \emph{$d(p,q)\le\delta+m$.} By \cref{lem:dc-ch7-2-8-normal-ball-dist}(1),
  $d(p,x_0)=|z_0|_p=\delta$, so the triangle inequality gives
  $d(p,q)\le d(p,x_0)+d(x_0,q)=\delta+m$.

  \emph{$d(p,q)\ge\delta+m$.} Suppose not; then there is a $C^1$ curve
  $\sigma:[0,1]\to M$ from $p$ to $q$ with $\ell(\sigma)<\delta+m$. First,
  $q$ lies outside the open region
  $U_\delta=\exp_p\{z\in B_\varepsilon(0):|z|_p<\delta\}$: a point
  $\exp_p v\in U_\delta$ has, by the Gram bound, $\|v\|\le\sqrt
  c\,|v|_p<\sqrt c\,\delta<\varepsilon$, hence
  $d(p,\exp_p v)=|v|_p<\delta$ by
  \cref{lem:dc-ch7-2-8-normal-ball-dist}(1), while $d(p,q)\ge\delta$.
  Since $\sigma(0)=p\in U_\delta$ and $\sigma(1)=q\notin U_\delta$, let
  $T=\inf\{t:\sigma(t)\notin U_\delta\}$ be the first exit time. For $t<T$
  the curve stays in $U_\delta$, whose closure is contained in the compact
  set $\exp_p(\{|z|_p\le\delta\}\cap\overline B_{\sqrt c\,\delta}(0))$; by
  continuity $\sigma(T)=\exp_p z$ with $|z|_p\le\delta$, and $|z|_p<\delta$
  would put $\sigma(T)$ back in the open set $U_\delta$, contradicting the
  definition of $T$; so $|z|_p=\delta$ exactly and $\sigma(T)\in
  S_\delta(p)$. On $[0,T]$ the curve stays in the normal region, so the
  radius comparison through the polar lift
  (\cref{lem:dc-ch7-2-8-normal-ball-dist}, competitor bound (i), applied on
  $[0,T]$ with \cref{lem:dc-ch7-2-8-chart-length}) gives
  $\ell(\sigma|_{[0,T]})\ge|z|_p=\delta$. Splitting the length at $T$ and
  bounding the tail by the distance,
  $$
  \ell(\sigma)\;=\;\ell(\sigma|_{[0,T]})+\ell(\sigma|_{[T,1]})
  \;\ge\;\delta+d(\sigma(T),q)\;\ge\;\delta+m,
  $$
  contradicting $\ell(\sigma)<\delta+m$. Taking the infimum over $\sigma$
  yields $d(p,q)\ge\delta+m$, hence equality. $\square$
\end{proof}

\begin{lemma}[the geodesic-sphere step]
  \label{lem:dc-ch7-2-8-step}
  \dcref{ch7:2.8}
  \uses{def:dc-ch7-2-4, lem:dc-ch7-2-8-ray-geodesic, lem:dc-ch7-2-8-sphere-min}
  \lean{Riemannian.Exponential.exists_minimizing_step}
  \leanok
  Let $M$ be a Riemannian manifold whose distance $d$ is the Riemannian
  distance of \cref{def:dc-ch7-2-4} and let $p\in M$. There is $\delta_0>0$
  such that for every $0<\delta<\delta_0$ and every $q\in M$ with
  $d(p,q)\ge\delta$ there is a continuous geodesic segment
  $\gamma:[0,1]\to M$ with
  $$
  \gamma(0)=p,\qquad d\big(p,\gamma(1)\big)=\delta,\qquad
  d(p,q)\;=\;\delta+d\big(\gamma(1),q\big).
  $$
  This is one full step of do Carmo's growth argument in the proof of
  \cref{thm:dc-ch7-2-8}: ride a radial geodesic to the point of the geodesic
  sphere $S_\delta(p)$ closest to $q$; the remaining distance drops by exactly
  $\delta$.
\end{lemma}
\begin{proof}
  \leanok
  \uses{lem:dc-ch7-2-8-ray-geodesic, lem:dc-ch7-2-8-sphere-min}
  Let $\varepsilon,c$ be the constants of \cref{lem:dc-ch7-2-8-sphere-min} at
  $p$ and $\rho,b$ those of \cref{lem:dc-ch7-2-8-ray-geodesic}; put
  $\delta_0=\min(\varepsilon,\rho)/\sqrt c$. For $0<\delta<\delta_0$ and
  $d(p,q)\ge\delta$, \cref{lem:dc-ch7-2-8-sphere-min} yields
  $z_0\in T_pM$ with $|z_0|_p=\delta$, $\|z_0\|\le\sqrt c\,\delta<\rho$,
  $d(p,\exp_p z_0)=\delta$ and $d(p,q)=\delta+d(\exp_p z_0,q)$. By
  \cref{lem:dc-ch7-2-8-ray-geodesic} the radial curve
  $\gamma(t)=\exp_p(t\,z_0)$ is a continuous geodesic on $[0,1]$ (indeed on
  $(-b,b)\supset[0,1]$) from $\gamma(0)=p$ to $\gamma(1)=\exp_p z_0$, which is
  exactly the displayed step. $\square$
\end{proof}

\begin{corollary}[compact $\Rightarrow$ complete]
  \label{cor:dc-ch7-2-9}
  \dcref{ch7:2.9}
  \uses{def:dc-ch7-2-2}
  \lean{Riemannian.Geodesic.isGeodesicallyComplete_of_compactSpace}
  \leanok
  If $M$ is compact then $M$ is complete.
\end{corollary}
\begin{proof}
  \leanok
  \uses{lem:dc-ch7-2-8-c-to-d}
  A compact metric space is metrically complete: a Cauchy sequence has a
  convergent subsequence, and a Cauchy sequence with a convergent subsequence
  converges. By \cref{thm:dc-ch7-2-8}, c) $\Rightarrow$ d)
  (\cref{lem:dc-ch7-2-8-c-to-d}), $M$ is geodesically complete. $\square$
\end{proof}

\begin{lemma}[a metric-preserving inclusion is $1$-Lipschitz for the Riemannian distances]
  \label{lem:dc-ch7-2-10-lipschitz}
  \dcref{ch7:2.10}
  \uses{def:dc-ch7-2-4, def:dc-ch1-2-2}
  \lean{Riemannian.DCPreservesMetric.enorm_mfderiv_eq,
    Riemannian.DCPreservesMetric.pathELength_comp_eq,
    Riemannian.DCPreservesMetric.riemannianEDist_le}
  \leanok
  Let $\iota:N\to M$ be a smooth map which \emph{preserves the metric} — a local isometry
  for the induced metric $g_N=\iota^*g_M$ (\cref{def:dc-ch1-2-2}), i.e.
  $\langle d\iota_p u,d\iota_p v\rangle_M=\langle u,v\rangle_N$ for all $p\in N$ and all
  $u,v\in T_pN$. Then $\iota$ does not increase the Riemannian distance:
  $d_M(\iota a,\iota b)\le d_N(a,b)$ for all $a,b\in N$.
\end{lemma}
\begin{proof}
  \leanok
  \uses{def:dc-ch7-2-4}
  Metric preservation gives $|d\iota_p(v)|=|v|$ on every tangent space, hence equal
  arclength integrands, so $\ell(\iota\circ c)=\ell(c)$ for every $C^1$ curve $c$. For any
  $C^1$ curve $c$ in $N$ from $a$ to $b$, the image $\iota\circ c$ is a $C^1$ curve in $M$
  from $\iota a$ to $\iota b$, whence $d_M(\iota a,\iota b)\le\ell(\iota\circ c)=\ell(c)$.
  Taking the infimum over all such $c$ (\cref{def:dc-ch7-2-4}) yields
  $d_M(\iota a,\iota b)\le d_N(a,b)$. (The reverse inequality can fail: a minimizing
  curve of $M$ between $\iota a$ and $\iota b$ need not lie in the submanifold.) $\square$
\end{proof}

\begin{corollary}[closed submanifold complete]
  \label{cor:dc-ch7-2-10}
  \dcref{ch7:2.10}
  \uses{lem:dc-ch7-2-10-lipschitz, prop:dc-ch7-2-6, thm:dc-ch7-2-8}
  \lean{Riemannian.completeSpace_of_isClosedEmbedding_dcPreservesMetric}
  \leanok
  A closed submanifold of a complete Riemannian manifold is complete in the induced
  metric; in particular, the closed submanifolds of Euclidean space are complete.
\end{corollary}
\begin{proof}
  \leanok
  \uses{lem:dc-ch7-2-10-lipschitz, prop:dc-ch7-2-6}
  Formalize the closed submanifold with its induced metric as a smooth \emph{closed
  embedding} $\iota:N\to M$ — a homeomorphism onto its closed range — which \emph{preserves}
  the induced metric $g_N=\iota^*g_M$, and let $M$ be complete. By
  \cref{lem:dc-ch7-2-10-lipschitz} $\iota$ is $1$-Lipschitz for the Riemannian distances.
  Let $(x_n)$ be a Cauchy sequence in $N$. Then $(\iota x_n)$ is Cauchy in the complete
  space $M$, so it converges to some $y\in M$. Since $\iota$ has closed range, $y=\iota x$
  for some $x\in N$; since $\iota$ is a homeomorphism onto its range, $\iota x_n\to\iota x$
  forces $x_n\to x$ in the manifold topology of $N$, which coincides with its metric
  topology (\cref{prop:dc-ch7-2-6}). Hence every Cauchy sequence of $N$ converges, i.e.
  $N$ is complete. Taking $M=\mathbb R^n$ shows that closed subsets of Euclidean space are
  complete. $\square$
\end{proof}

\section{The Theorem of Hadamard}

As an application of the theorem of Hopf-Rinow, we are going to prove the
following global fact.

\begin{theorem}[Hadamard]
  \label{thm:dc-ch7-3-1}
  \dcref{ch7:3.1}
  \lean{Riemannian.Jacobi.hadamardDiffeomorphOfNonpos_complete, Riemannian.Jacobi.hadamardDiffeomorphOfNonpos_complete_coe}
  \uses{lem:dc-ch7-3-2, lem:dc-ch7-3-3, lem:dc-ch7-3-1-covering-diffeo, lem:dc-ch7-3-1-assembly, lem:dc-ch7-3-4-pullback-metric, lem:dc-ch7-3-4-pole-diffeo, lem:dc-ch7-3-expmap-global, lem:dc-ch7-3-expmap-smooth, thm:dc-ch7-2-8}
  \leanok
  Let $M$ be a complete Riemannian manifold, simply connected, with sectional
  curvature $K(p,\sigma)\le 0$ for all $p\in M$ and for all $\sigma\subset T_p(M)$.
  Then $M$ is diffeomorphic to $\mathbb{R}^n$, $n=\dim M$; more precisely,
  $\exp_p:T_pM\to M$ is a diffeomorphism.

  Before starting the proof we need a few lemmas. The following lemma shows that the
  exponential map of a manifold with non-positive curvature is a local
  diffeomorphism.

  The theorem is fully assembled (\lstinline{hadamardDiffeomorphOfNonpos_complete} produces a
  \lstinline{Diffeomorph}~$T_pM\simeq M$, and \lstinline{hadamardDiffeomorphOfNonpos_complete_coe}
  guards that this diffeomorphism \emph{is} $\exp_p$) from the two halves: the differential of
  $\exp_p$ is a linear isomorphism under $K\le 0$ (\cref{lem:dc-ch7-3-2}), and the
  poles/covering upgrade of a smooth local diffeomorphism to a diffeomorphism
  (\cref{lem:dc-ch7-3-4-pole-diffeo}), together with the global $C^\infty$ smoothness of
  $\exp_p$ (\cref{lem:dc-ch7-3-expmap-smooth}), which was the last remaining input and is now
  proved. The proof is that of do Carmo: \cref{lem:dc-ch7-3-2} makes $\exp_p$ a local
  diffeomorphism, and since $M$ is complete the pulled-back metric $\exp_p^{*}g$ makes
  $\exp_p$ a metric-expanding local diffeomorphism out of a (geodesically) complete manifold,
  hence a covering map (\cref{lem:dc-ch7-3-3}); as $M$ is simply connected this covering is a
  diffeomorphism.
\end{theorem}

\begin{lemma}[the Jacobi energy is convex, and $J$ has no interior zero, when $K\le 0$]
  \label{lem:dc-ch7-3-2-no-conjugate}
  \dcref{ch7:3.2}
  \lean{Riemannian.Jacobi.jacobiCoefOp_quadratic_nonpos,
    Riemannian.Jacobi.frameJacobi_ne_zero_of_nonpos}
  \leanok
  \uses{def:dc-ch5-2-1, lem:dc-ch5-2-1-real-curvature}
  Let $J=\sum_i f_i e_i$ be a Jacobi field in a parallel orthonormal frame along the
  geodesic $\gamma$, with $J(0)=0$ and nonzero initial velocity $J'(0)\ne 0$. Under
  nonpositive sectional curvature, read in the chart as $\langle R(\gamma',w)\gamma',w\rangle\le 0$
  for every $w$, the energy $q=\langle J,J\rangle$ satisfies
  $$
  \langle J,J\rangle''=2|J'|^2-2\langle R(\gamma',J)\gamma',J\rangle\ge 0,
  $$
  because the curvature term $\langle R(\gamma',J)\gamma',J\rangle=\sum_{ij}a_{ij}f_if_j\le 0$
  is the (negative semidefinite) quadratic form of the Jacobi coefficient
  $a_{ij}=\langle R(\gamma',e_i)\gamma',e_j\rangle$. Since $q(0)=0$ and $q'(0)=0$, convexity
  gives $q'\ge 0$, so $q$ is nondecreasing; nontriviality ($J'(0)\ne 0$) then forces
  $q(t)>0$ for every $t>0$, i.e. $J(t)\ne 0$. Hence no point $\gamma(t)$, $t>0$, is
  conjugate to $\gamma(0)$.
\end{lemma}

\begin{lemma}[no conjugate points when $K\le 0$]
  \label{lem:dc-ch7-3-2}
  \dcref{ch7:3.2}
  \lean{Riemannian.Jacobi.not_isConjugatePointAt_one_of_nonpos,
    Riemannian.Jacobi.expDifferential_isEquiv_of_nonpos,
    Riemannian.Jacobi.expMapGlobal_locallyInjective_of_nonpos,
    Riemannian.Jacobi.isLocalDiffeomorph_expMapGlobal_hadamard_of_nonpos_complete}
  \leanok
  \uses{def:dc-ch5-2-1, def:dc-ch5-3-1, def:dc-ch5-3-4, lem:dc-ch7-3-2-no-conjugate, lem:dc-ch7-3-expmap-smooth}
  Let $M$ be a complete Riemannian manifold with $K(p,\sigma)\le 0$ for all $p\in M$
  and for all $\sigma\subset T_pM$. Then for all $p\in M$ the conjugate locus
  $C(p)=\varnothing$; in particular the exponential map $\exp_p:T_pM\to M$ is a local
  diffeomorphism.

  \emph{Formalization note.} The nonpositive-curvature hypothesis is carried in the operator
  (equivalently chart) form $\langle R(a,c)c,a\rangle_g\ge 0$, i.e.
  $\langle R(\gamma',J)\gamma',J\rangle\le 0$ — the same convention as
  \cref{lem:dc-ch7-3-2-no-conjugate}, and by skew-adjointness of the curvature operator this is
  do Carmo's $K\le 0$. The three ingredients are all axiom-clean: $C(p)=\varnothing$
  (\lstinline{not_isConjugatePointAt_one_of_nonpos}, via the intrinsic energy argument below run
  along the whole geodesic, patched across charts); $d(\exp_p)_v$ a linear isomorphism at every
  $v$ (\lstinline{expDifferential_isEquiv_of_nonpos}); and $\exp_p$ injective near every $v$
  (\lstinline{expMapGlobal_locallyInjective_of_nonpos}, inverse function theorem on the chart
  reading). With the global $C^\infty$ smoothness of $\exp_p$ now proved
  (\cref{lem:dc-ch7-3-expmap-smooth}), the full $C^\infty$ local-diffeomorphism packaging
  \lstinline{isLocalDiffeomorph_expMapGlobal_hadamard_of_nonpos_complete} is unconditional: at
  each $v$ the chart-reading strict-derivative isomorphism, through the manifold inverse function
  theorem and composed with the chart inverse, gives $\exp_p$ as a $C^\infty$ local
  diffeomorphism, which \cref{thm:dc-ch7-3-1} feeds into the covering-space assembly.
\end{lemma}
\begin{proof}
  \uses{lem:dc-ch7-3-2-no-conjugate, prop:dc-ch5-3-5}
  Let $J$ be a non-trivial (that is, not identically zero) Jacobi field
  along a geodesic $\gamma:[0,\infty)\to M$, where $\gamma(0)=p$ and $J(0)=0$. Then
  from the hypothesis on the curvature and from the Jacobi equation
  
  $$
  \begin{aligned}
  \langle J,J\rangle''&=2\langle J',J'\rangle+2\langle J'',J\rangle\\
  &=2\langle J',J'\rangle-2\langle R(\gamma',J)\gamma',J\rangle\\
  &=2|J'|^2-2K(\gamma',J)|\gamma'\wedge J|^2\ge 0.
  \end{aligned}
  $$
  
  Therefore $\langle J,J\rangle'(t_2)\ge\langle J,J\rangle'(t_1)$ whenever
  $t_2>t_1$. Since $J'(0)\ne 0$ and $\langle J,J\rangle'(0)=0$, it follows in
  addition that for $t$ a sufficiently small positive number
  
  $$
  \langle J,J\rangle(t)>\langle J,J\rangle(0).
  $$
  
  It follows that for all $t>0$, $\langle J,J\rangle(t)>0$, and $\gamma(t)$ is not
  conjugate to $\gamma(0)$ along $\gamma$. $\square$
  
  The crucial point in the proof of the Hadamard theorem is given in the lemma
  below, which is of independent interest.
\end{proof}

\begin{lemma}[length expansion under a metric-expanding map]
  \label{lem:dc-ch7-3-3-length-expansion}
  \dcref{ch7:3.3}
  \uses{def:dc-ch7-2-4, def:dc-ch1-2-9}
  \lean{Riemannian.DCExpandsMetric, Riemannian.DCExpandsMetric.pathELength_le, Riemannian.DCExpandsMetric.riemannianEDist_le_pathELength_comp}
  \leanok
  Let $f:M\to N$ be a differentiable map between Riemannian manifolds which
  \emph{expands the metric}: $|df_p(v)|\ge|v|$ for every $p\in M$ and every
  $v\in T_pM$. Then for every piecewise differentiable curve
  $\bar c:[a,b]\to M$,
  $$
  \ell(f\circ\bar c)\ \ge\ \ell(\bar c)\ \ge\ d\big(\bar c(a),\bar c(b)\big).
  $$
  In particular $d\big(\bar c(a),\bar c(b)\big)\le\ell(f\circ\bar c)$: the
  source distance between the endpoints of a curve is bounded by the length of
  its image. This is the displayed estimate in the proof of
  \cref{lem:dc-ch7-3-3} that keeps the lifted points inside a bounded set.
\end{lemma}
\begin{proof}
  \leanok
  \uses{def:dc-ch7-2-4}
  By the chain rule the velocity of $f\circ\bar c$ is
  $(f\circ\bar c)'(t)=df_{\bar c(t)}\big(\bar c'(t)\big)$, so the expansion
  hypothesis gives
  $|(f\circ\bar c)'(t)|=\big|df_{\bar c(t)}(\bar c'(t))\big|\ge|\bar c'(t)|$ at
  every $t$. Integrating this pointwise inequality over $[a,b]$ (both lengths
  are the tangent integral $\int_a^b|\,\cdot\,'|$, and the integral is
  monotone) gives
  $\ell(f\circ\bar c)=\int_a^b|(f\circ\bar c)'|\ge\int_a^b|\bar c'|=\ell(\bar c)$.
  Finally $\bar c$ is itself a piecewise differentiable curve joining
  $\bar c(a)$ to $\bar c(b)$, so its length is at least the infimum
  $d\big(\bar c(a),\bar c(b)\big)$ over all such curves, by
  \cref{def:dc-ch7-2-4}. $\square$
\end{proof}

\begin{lemma}[lifts of finite-length curves are relatively compact]
  \label{lem:dc-ch7-3-3-relatively-compact-lift}
  \dcref{ch7:3.3}
  \uses{lem:dc-ch7-3-3-length-expansion, thm:dc-ch7-2-8, def:dc-ch7-2-4}
  \lean{Riemannian.DCExpandsMetric.dist_le_pathELength_comp,
    Riemannian.DCExpandsMetric.lift_mem_closedBall,
    Riemannian.DCExpandsMetric.lift_image_subset_isCompact}
  \leanok
  Let $f:M\to N$ expand the metric ($|df_p(v)|\ge|v|$ for all $p,v$), and let $M$
  be a complete Riemannian manifold. Then for every piecewise differentiable
  curve $\bar c:[a,b]\to M$ whose image $f\circ\bar c$ has finite length, and for
  every $t\in[a,b]$,
  $$
  d\big(\bar c(a),\bar c(t)\big)\ \le\ \ell(f\circ\bar c),
  $$
  so $\bar c(t)$ lies in the closed metric ball $\bar B\big(\bar c(a),\ell(f\circ\bar c)\big)$.
  Consequently the whole image $\bar c([a,b])$ is contained in a single
  \emph{compact} subset $K\subset M$.
\end{lemma}
\begin{proof}
  \leanok
  \uses{lem:dc-ch7-3-3-length-expansion, thm:dc-ch7-2-8}
  Fix $t\in[a,b]$. The restriction $\bar c|_{[a,t]}$ is a piecewise
  differentiable curve, so by the length-expansion estimate
  (\cref{lem:dc-ch7-3-3-length-expansion}) applied on $[a,t]$ and monotonicity of
  length in the right endpoint,
  $$
  d\big(\bar c(a),\bar c(t)\big)\ \le\ \ell\big(f\circ\bar c|_{[a,t]}\big)\ \le\ \ell(f\circ\bar c)=:R,
  $$
  and $R<\infty$ by hypothesis. Hence $\bar c(t)\in\bar B(\bar c(a),R)$ for every
  $t$, i.e. $\bar c([a,b])\subseteq\bar B(\bar c(a),R)$. Since $M$ is complete, by
  the Hopf–Rinow theorem (\cref{thm:dc-ch7-2-8}, c) $\Rightarrow$ b)) every closed
  bounded subset of $M$ is compact; in particular $K:=\bar B(\bar c(a),R)$ is
  compact. $\square$
\end{proof}

\begin{lemma}[local liftability of a curve through a local homeomorphism]
  \label{lem:dc-ch7-3-3-local-lift}
  \dcref{ch7:3.3}
  \lean{Riemannian.IsLocalHomeomorph.exists_rightward_lift,
    Riemannian.IsLocalHomeomorph.exists_extend_lift}
  \leanok
  Let $f:M\to N$ be a local homeomorphism and $c:[0,1]\to N$ a continuous curve.
  \begin{enumerate}
    \item For every $s$ and every point $q\in M$ over $c(s)$ (that is,
      $f(q)=c(s)$) the curve $c$ lifts through $q$ on a nondegenerate interval to
      the right: there is a continuous $\bar c:[s,s+\varepsilon]\to M$ with
      $\bar c(s)=q$ and $f\circ\bar c=c$.
    \item Consequently, if $c$ lifts to a continuous $\bar c:[a,t]\to M$ with
      $f\circ\bar c=c$, then the lift extends to $[a,t+\delta]$ for some
      $\delta>0$: the set of parameters over which $c$ lifts starting from a fixed
      initial point is open to the right.
  \end{enumerate}
\end{lemma}
\begin{proof}
  \leanok
  For (1), choose an open partial homeomorphism $\varphi$ of $f$ with $q$ in its
  source and $f=\varphi$ there. Then $c(s)=f(q)=\varphi(q)$ lies in the open set
  $\varphi(\text{source})$, so $c^{-1}(\varphi(\text{source}))$ is an open
  neighbourhood of $s$ and contains some $[s,s+\varepsilon]$; on it
  $\bar c:=\varphi^{-1}\circ c$ is continuous, satisfies $\bar c(s)=q$ (as
  $\varphi^{-1}(\varphi(q))=q$) and $f\circ\bar c=c$ (as
  $\varphi(\varphi^{-1}(y))=y$ on $\varphi$'s target).
  For (2), apply (1) at the endpoint value $q:=\bar c(t)$ to obtain a rightward
  lift $\bar c_1$ on $[t,t+\delta]$ with $\bar c_1(t)=\bar c(t)$; the two lifts
  agree at $t$, so their piecewise gluing is a continuous lift on
  $[a,t+\delta]$. $\square$
\end{proof}

\begin{lemma}[path lifting for a metric-expanding local diffeomorphism]
  \label{lem:dc-ch7-3-3-path-lifting}
  \dcref{ch7:3.3}
  \uses{thm:dc-ch7-2-8, def:dc-ch7-2-4, lem:dc-ch7-3-3-length-expansion, lem:dc-ch7-3-3-relatively-compact-lift, lem:dc-ch7-3-3-local-lift}
  \lean{Riemannian.DCExpandsMetric.exists_pathLift, Riemannian.IsLocalDiffeomorph.contMDiffOn_lift, Riemannian.DCExpandsMetric.dist_lift_le}
  \leanok
  Let $M$ be a complete Riemannian manifold and let $f:M\to N$ be a local
  diffeomorphism which expands the metric ($|df_p(v)|\ge|v|$ for all $p\in M$ and
  all $v\in T_pM$). Then $f$ has the \emph{path-lifting property}: for every
  $C^1$ curve $c:[0,1]\to N$ of finite length and every $q\in M$ with $f(q)=c(0)$,
  there exists a continuous curve $\bar c:[0,1]\to M$ with $\bar c(0)=q$ and
  $f\circ\bar c=c$.

  Moreover any continuous lift of a $C^1$ curve through a local diffeomorphism is
  itself $C^1$ (locally it equals $(f|_V)^{-1}\circ c$ for a local inverse), and
  the endpoints $\bar c(t)$ of a partial lift satisfy
  $d\big(\bar c(0),\bar c(t)\big)\le\ell(c)$ (the estimate that traps the lift in a
  compact ball).
\end{lemma}
\begin{proof}
  \leanok
  \uses{thm:dc-ch7-2-8, lem:dc-ch7-3-3-length-expansion, lem:dc-ch7-3-3-relatively-compact-lift, lem:dc-ch7-3-3-local-lift}
  Let $A\subseteq[0,1]$ be the set of parameters $t$ such that $c|_{[0,t]}$ admits
  a continuous lift starting at $q$. Then $0\in A$ (the constant curve at $q$
  lifts $c|_{\{0\}}$), and $A$ is downward closed, since restricting a lift on
  $[0,t]$ to $[0,t']$ with $t'\le t$ lifts $c|_{[0,t']}$. Put $T=\sup A$; then
  $T\in[0,1]$, and every $t$ with $0\le t<T$ lies in $A$.

  \emph{$A$ is closed, i.e. $T\in A$.} We may assume $T>0$ (otherwise $T=0\in A$).
  For $0\le t<T$ choose a lift $\bar c_t$ of $c|_{[0,t]}$ from $q$; being a
  continuous lift of a $C^1$ curve through the local diffeomorphism $f$, it is
  $C^1$, so by the length estimate (\cref{lem:dc-ch7-3-3-relatively-compact-lift},
  which combines \cref{lem:dc-ch7-3-3-length-expansion} with the properness of the
  complete manifold $M$, \cref{thm:dc-ch7-2-8} c)$\Rightarrow$b)) its endpoint
  $\bar c_t(t)$ lies in the fixed compact ball $K=\overline B\big(q,\ell(c)\big)$.
  Hence as $t\to T^-$ the endpoints $\bar c_t(t)$ have an accumulation point
  $r\in K$. Since $f$ is continuous and $f\big(\bar c_t(t)\big)=c(t)\to c(T)$, we
  get $f(r)=c(T)$. Let $V$ be a neighbourhood of $r$ on which $f$ restricts to a
  homeomorphism; then $c(T)=f(r)\in f(V)$, so by continuity $c(t)\in f(V)$ for all
  $t$ in a left neighbourhood $(T-\eta,T]$ of $T$, and $\bar c_{t_*}(t_*)\in V$ for
  some $t_*\in(T-\eta,T)$. On $(T-\eta,T]$ the map $g=(f|_V)^{-1}\circ c$ is a
  continuous lift of $c$ with $g(t_*)=\bar c_{t_*}(t_*)$ (both project to $c(t_*)$
  and $f|_V$ is injective). Gluing $\bar c_{t_*}$ on $[0,t_*]$ to $g$ on $[t_*,T]$
  yields a continuous lift of $c|_{[0,T]}$ from $q$, so $T\in A$.

  \emph{$A$ is open to the right.} If $t\in A$ and $t<1$, the openness ingredient
  (\cref{lem:dc-ch7-3-3-local-lift}) extends the lift of $c|_{[0,t]}$ to
  $[0,t+\delta]$ for some $\delta>0$, so $\min(t+\delta,1)\in A$.

  Since $T\in A$, if $T<1$ the previous paragraph would place a parameter
  $>T$ in $A$, contradicting $T=\sup A$. Hence $T=1\in A$, and $c$ lifts on all of
  $[0,1]$. $\square$
\end{proof}

\begin{lemma}[uniqueness of continuous lifts]
  \label{lem:dc-ch7-3-3-lift-unique}
  \dcref{ch7:3.3}
  \uses{lem:dc-ch7-3-3-path-lifting}
  \lean{Riemannian.IsLocalHomeomorph.eqOn_lift, Riemannian.IsLocalDiffeomorph.eqOn_pathLift, Riemannian.DCExpandsMetric.existsUnique_pathLift_eqOn}
  \leanok
  Let $f:M\to N$ be a local diffeomorphism out of a Riemannian manifold $M$ (in
  particular a Hausdorff space). If $\bar c_1,\bar c_2:S\to M$ are two continuous
  lifts of the same curve $c$ (i.e.\ $f\circ\bar c_i=c$) over a connected parameter
  set $S$, and $\bar c_1(t_0)=\bar c_2(t_0)$ for some $t_0\in S$, then
  $\bar c_1=\bar c_2$ on all of $S$. In particular the continuous lift of a $C^1$
  curve $c:[0,1]\to N$ through a chosen point $q$ over $c(0)$, produced by
  \cref{lem:dc-ch7-3-3-path-lifting}, is \emph{unique}.
\end{lemma}
\begin{proof}
  \leanok
  \uses{lem:dc-ch7-3-3-path-lifting}
  A local diffeomorphism is a local homeomorphism, hence \emph{locally injective}:
  each point of $M$ has an open neighbourhood on which $f$ is injective. Since $M$
  is a metric space it is Hausdorff, so $f$ is a \emph{separated map} (any two
  distinct points of $M$ have disjoint neighbourhoods, a fortiori two distinct
  points of the same fibre). For a separated, locally injective map the
  \emph{equal locus} $\{t\in S:\bar c_1(t)=\bar c_2(t)\}$ of two continuous lifts is
  both closed (separatedness) and open (local injectivity), hence clopen; as it is
  nonempty ($t_0$ lies in it) and $S$ is connected, it is all of $S$. Applying this
  to $S=[0,1]$ (connected) and the two lifts starting at $q$ gives uniqueness of the
  lift of \cref{lem:dc-ch7-3-3-path-lifting}. $\square$
\end{proof}

\begin{lemma}[finite length of a chart-contained $C^1$ curve]
  \label{lem:dc-ch7-3-3-chart-length-finite}
  \dcref{ch7:3.3}
  \uses{def:dc-ch1-2-9}
  \lean{Riemannian.Geodesic.pathELength_ne_top_of_mapsTo_chartSource}
  \leanok
  If $\gamma:[a,b]\to N$ is a $C^1$ curve whose image lies in the source of a single
  chart of $N$, then its length $\ell(\gamma)$ is finite. Reading the length through
  the chart presents it as $\mathrm{ofReal}\big(\int_a^b\sqrt{\langle\gamma',\gamma'\rangle}\big)$,
  a real number. This supplies the finite-length hypothesis of
  \cref{lem:dc-ch7-3-3-path-lifting} for the straight chart-segments used to build
  evenly covered neighbourhoods.
\end{lemma}
\begin{proof}
  \leanok
  \uses{def:dc-ch1-2-9}
  The chart reading of the length (do Carmo's $\ell(\gamma)=\int|\gamma'|$, computed
  in chart coordinates) equals $\mathrm{ofReal}$ of a Riemann integral of the
  continuous integrand $\sqrt{\langle\gamma',\gamma'\rangle}$ over the compact
  $[a,b]$; an $\mathrm{ofReal}$ of a real number is never $+\infty$. $\square$
\end{proof}

\begin{lemma}[closed image; surjectivity onto a connected target]
  \label{lem:dc-ch7-3-3-closed-image}
  \dcref{ch7:3.3}
  \uses{lem:dc-ch7-3-3-path-lifting, lem:dc-ch7-3-3-chart-length-finite}
  \lean{Riemannian.DCExpandsMetric.isClosed_range, Riemannian.DCExpandsMetric.surjective}
  \leanok
  Let $M$ be a complete Riemannian manifold and $f:M\to N$ a local diffeomorphism
  which expands the metric. Then the image $f(M)$ is \emph{closed} in $N$; if $N$ is
  connected and $M$ nonempty, $f$ is \emph{surjective}. This handles the points of
  $N$ outside the image in the covering-map conclusion: at such a point the fibre is
  empty, and an evenly covered neighbourhood is one disjoint from the image — which
  exists precisely because the image is closed.
\end{lemma}
\begin{proof}
  \leanok
  \uses{lem:dc-ch7-3-3-path-lifting, lem:dc-ch7-3-3-chart-length-finite}
  The image is open, since $f$ is a local diffeomorphism (an open map). For
  closedness, let $y$ be a limit point of $f(M)$. Work in a chart at $y$ and pick a
  ball $B$ around $y$ (in chart coordinates) inside the chart target; some image
  point $f(x_1)$ lies in $B$. The straight chart-segment from $f(x_1)$ to $y$ is a
  $C^1$ curve confined to one chart, hence of finite length
  (\cref{lem:dc-ch7-3-3-chart-length-finite}), so it lifts through $x_1$
  (\cref{lem:dc-ch7-3-3-path-lifting}); the lift's endpoint maps to $y$, so
  $y\in f(M)$. Thus $f(M)$ is clopen; a nonempty clopen subset of a connected $N$ is
  all of $N$. $\square$
\end{proof}

\begin{lemma}[expansion $\Rightarrow$ covering map]
  \label{lem:dc-ch7-3-3}
  \dcref{ch7:3.3}
  \uses{thm:dc-ch7-2-8, lem:dc-ch7-3-3-length-expansion, lem:dc-ch7-3-3-relatively-compact-lift, lem:dc-ch7-3-3-local-lift, lem:dc-ch7-3-3-path-lifting, lem:dc-ch7-3-3-lift-unique, lem:dc-ch7-3-3-chart-length-finite, lem:dc-ch7-3-3-closed-image}
  \lean{Riemannian.DCExpandsMetric.isCoveringMap, Riemannian.DCExpandsMetric.exists_trivialization_at, Riemannian.IsLocalHomeomorph.discreteTopology_fiber}
  \leanok
  Let $M$ be a complete Riemannian manifold and let $f:M\to N$ be a local
  diffeomorphism onto a Riemannian manifold $N$ which has the following property:
  for all $p\in M$ and for all $v\in T_pM$, we have $|df_p(v)|\ge|v|$. Then $f$ is a
  covering map.
\end{lemma}
\begin{proof}
  \leanok
  By a general property of covering spaces (cf. M. do Carmo, [dC 2],
  p. 383), it suffices to show that $f$ has the path lifting property for curves in
  $N$: given a differentiable curve $c:[0,1]\to N$ and a point $q\in M$ with
  $f(q)=c(0)$, there exists a curve $\bar c:[0,1]\to M$ with $\bar c(0)=q$ and
  $f\circ\bar c=c$. This path-lifting property is established in
  \cref{lem:dc-ch7-3-3-path-lifting}, and its lift is \emph{unique}
  (\cref{lem:dc-ch7-3-3-lift-unique}). We convert path lifting into the covering-map
  structure (evenly covered neighbourhoods) directly, by the following concrete
  construction.

  Fix a point $y_0\in N$ and a chart $\psi$ at $y_0$; let $V$ be a $\psi$-ball around
  $y_0$. Each point $y\in V$ is joined to $y_0$ by the $C^1$ segment
  $t\mapsto\psi^{-1}\big((1-t)\psi(y_0)+t\,\psi(y)\big)$ (staying in one chart, hence
  of finite length by \cref{lem:dc-ch7-3-3-chart-length-finite}). For each $e$ in the
  fibre $f^{-1}(y_0)$, lifting this family of segments through $e$ produces — by joint
  continuity of lifts of a continuous family of paths through a separated local
  homeomorphism (mathlib \texttt{IsLocalHomeomorph.continuous\_lift}, fed by
  \cref{lem:dc-ch7-3-3-path-lifting}) — a \emph{continuous section}
  $\varphi_e:V\to M$ with $f\circ\varphi_e=\mathrm{id}_V$ and $\varphi_e(y_0)=e$. A
  continuous section of a local homeomorphism is an open embedding, so each sheet
  $\varphi_e(V)$ is open; uniqueness of lifts (\cref{lem:dc-ch7-3-3-lift-unique})
  makes the sheets pairwise disjoint. Given any $m\in f^{-1}(V)$, lifting the reverse
  segment $y=f(m)\rightsquigarrow y_0$ from $m$ lands on a fibre point $e$, and by
  uniqueness the reversed lift is $\varphi_e$ along the forward segment, so
  $\varphi_e(y)=m$: the sheets \emph{exhaust} $f^{-1}(V)=\bigsqcup_e\varphi_e(V)$,
  which is exactly an evenly covered neighbourhood (packaged as a
  \texttt{Trivialization} over the discrete fibre $f^{-1}(y_0)$ — discrete because $f$
  is locally injective). The fibre being discrete, this feeds mathlib's
  \texttt{IsCoveringMap.mk'}; the points of $N$ \emph{outside} the image are handled
  by its closed-range branch, since the image is closed
  (\cref{lem:dc-ch7-3-3-closed-image}) so each such point has a neighbourhood with
  empty preimage. This completes the covering-map conclusion.

  For completeness we recall do Carmo's original clopen argument for path lifting,
  formalised in \cref{lem:dc-ch7-3-3-path-lifting}. Since $f$ is a local
  diffeomorphism at $q$, there exists $\varepsilon>0$ such that
  it is possible to define $\bar c:[0,\varepsilon]\to M$ with $\bar c(0)=q$ and
  $f\circ\bar c=c$; that is, $c$ can be lifted to a small interval starting from $q$.
  Because $f$ is a local diffeomorphism over all of $M$, the set of values
  $A\subset[0,1]$ such that $c$ can be lifted on $A$ starting from $q$ is an open
  interval on the right; that is, $A=[0,t_0)$. If we can show that $t_0\in A$, we
  shall have $A$ open and closed in $[0,1]$, therefore $A=[0,1]$ and $c$ can be
  lifted on the entire interval.
  
  To show that $t_0\in A$, let $\{t_n\}$, $n=1,\dots$, be an increasing sequence in
  $A$ with $\lim t_n=t_0$. Then the sequence $\{\bar c(t_n)\}$ is contained in a
  compact set $K\subset M$. Indeed, if this were not the case, since $M$ is
  complete, the distance from $\bar c(t_n)$ to $\bar c(0)$ would be arbitrarily
  large. However, by hypothesis,
  
  $$
  \ell_{0,t_n}(c)=\int_0^{t_n}\Big|\frac{dc}{dt}\Big|dt=\int_0^{t_n}\Big|df_{\bar c(t)}\Big(\frac{d\bar c}{dt}\Big)\Big|dt\ge\int_0^{t_n}\Big|\frac{d\bar c}{dt}\Big|dt\ge d(\bar c(t_n),\bar c(0)),
  $$
  
  implying that the length of $c$ between $0$ and $t_0$ is arbitrarily large, which
  is absurd, and proves the assertion.
  
  Since $\{\bar c(t_n)\}\subset K$, $n=1,\dots$, there exists an accumulation point
  $r\in M$ of $\{\bar c(t_n)\}$. Let $V$ be a neighborhood of $r$ such that $f\,|\,V$
  is a diffeomorphism. Then $c(t_0)\in f(V)$ and, by continuity, there exists an
  interval $I\subset[0,1]$, $t_0\in I$, such that $c(I)\subset f(V)$. Choose an
  index $n$ such that $\bar c(t_n)\in V$ and consider the lifting $g$ of $c$ on $I$
  passing through $r$. The liftings $g$ and $\bar c$ coincide on $[0,t_n)\cap I$,
  because $f\,|\,V$ is bijective. Therefore $g$ is an extension of $\bar c$ to $I$,
  hence $\bar c$ is defined at $t_0$ and $t_0\in A$. $\square$
\end{proof}

\begin{lemma}[a covering map onto a simply connected base is a homeomorphism]
  \label{lem:dc-ch7-3-1-covering-simplyconnected}
  \dcref{ch7:3.1}
  \lean{IsCoveringMap.bijective_of_simplyConnected, IsCoveringMap.homeomorphOfSimplyConnected}
  \leanok
  Let $f:E\to X$ be a covering map whose total space $E$ is connected and nonempty,
  and whose base $X$ is simply connected and locally path connected. Then $f$ is a
  homeomorphism; in particular it is bijective.
\end{lemma}
\begin{proof}
  \leanok
  Because $X$ is simply connected and locally path connected, the identity map of
  $X$ lifts through the covering map $f$ to a continuous global section $s:X\to E$,
  that is $f\circ s=\mathrm{id}_X$, uniquely determined by its value $s(x_0)=e_0$ at
  one chosen base point $x_0=f(e_0)$. The relation $f\circ s=\mathrm{id}_X$ exhibits
  $s$ as a right inverse of $f$, so $f$ is surjective. For injectivity, consider the
  two continuous self-maps $s\circ f$ and $\mathrm{id}_E$ of $E$: both are lifts of
  $f$ through $f$ (indeed $f\circ(s\circ f)=f=f\circ\mathrm{id}_E$), and they agree at
  $e_0$ (since $s(f(e_0))=s(x_0)=e_0$). A covering map is a separated, locally
  injective map, so any two continuous lifts of a fixed map over the connected space
  $E$ that agree at one point agree everywhere; hence $s\circ f=\mathrm{id}_E$ and $f$
  is injective. Finally a continuous open bijection is a homeomorphism, and a covering
  map is open. $\square$
\end{proof}

\begin{lemma}[metric-expanding local diffeomorphism onto a simply connected manifold]
  \label{lem:dc-ch7-3-1-covering-diffeo}
  \dcref{ch7:3.1}
  \uses{lem:dc-ch7-3-3, lem:dc-ch7-3-1-covering-simplyconnected}
  \lean{Riemannian.DCExpandsMetric.diffeomorphOfSimplyConnected}
  \leanok
  Let $f:M\to N$ be a smooth local diffeomorphism between Riemannian manifolds which
  expands the metric ($|df_p(v)|\ge|v|$ for all $p\in M$ and all $v\in T_pM$), let $M$
  be complete, and let $N$ be simply connected. Then $f$ is a diffeomorphism.
\end{lemma}
\begin{proof}
  \leanok
  \uses{lem:dc-ch7-3-3, lem:dc-ch7-3-1-covering-simplyconnected}
  By \cref{lem:dc-ch7-3-3}, $f$ is a covering map. Its total space $M$ is connected
  and its base $N$ is simply connected, so by
  \cref{lem:dc-ch7-3-1-covering-simplyconnected} the covering map $f$ is bijective.
  A bijective local diffeomorphism is a diffeomorphism, hence $f$ is a
  diffeomorphism. $\square$
\end{proof}

\begin{lemma}[the global exponential map of a complete manifold]
  \label{lem:dc-ch7-3-expmap-global}
  \dcref{ch7:2.2}
  \uses{def:dc-ch7-2-2, thm:dc-ch7-2-8}
  \lean{Riemannian.Geodesic.intrinsicMaximalGeodesic_eq_globalGeodesic, Riemannian.Geodesic.globalGeodesic, Riemannian.Geodesic.globalGeodesic_eq, Riemannian.Geodesic.globalGeodesic_smul, Riemannian.Exponential.expMapIntrinsic_eq_expMapGlobal, Riemannian.Exponential.expMapGlobal, Riemannian.Exponential.expMapGlobal_eq_of_isGeodesic, Riemannian.Exponential.expMapGlobal_smul, Riemannian.Exponential.expMapIntrinsic_surjective}
  \leanok
  Let $M$ be a complete Riemannian manifold and $p\in M$. For every $v\in T_pM$ there is a
  geodesic $\gamma_v:\mathbb{R}\to M$, defined on all of $\mathbb{R}$, with $\gamma_v(0)=p$
  and $\gamma_v'(0)=v$, and it is the unique such geodesic. The exponential map
  $\exp_p:T_pM\to M$, $\exp_p(v)=\gamma_v(1)$, is therefore defined on all of $T_pM$, and it
  is traced radially: $\exp_p(cv)=\gamma_v(c)$ for all $c\in\mathbb{R}$, so the curves
  $c\mapsto\exp_p(cv)$ are precisely the geodesics of $M$ issuing from $p$. If moreover $M$
  is connected, then $\exp_p$ is surjective.
\end{lemma}
\begin{proof}
  \leanok
  \uses{def:dc-ch7-2-2, thm:dc-ch7-2-8}
  Existence of a geodesic defined on all of $\mathbb{R}$ with prescribed initial data
  $(p,v)$ is the $(c)\Rightarrow(d)$ direction of \cref{thm:dc-ch7-2-8}. Two geodesics
  defined on all of $\mathbb{R}$ that agree in position and velocity at one time agree
  everywhere, by uniqueness of solutions of the geodesic equation on the connected parameter
  line $\mathbb{R}$; this makes $\gamma_v$, hence $\exp_p$, well defined independently of the
  chart in which the initial velocity is read. For radial homogeneity, the affine
  reparametrisation $s\mapsto\gamma_v(cs)$ is again a geodesic based at $p$ with initial
  velocity $cv$, so by uniqueness $\gamma_{cv}(s)=\gamma_v(cs)$; evaluating at $s=1$ gives
  $\exp_p(cv)=\gamma_v(c)$. Finally, when $M$ is connected, the length-minimizing-geodesic
  conclusion of \cref{thm:dc-ch7-2-8} joins every $q\in M$ to $p$ by a geodesic defined on
  all of $\mathbb{R}$; reading its initial velocity $v$ gives $q=\gamma_v(1)=\exp_p(v)$, so
  $\exp_p$ is surjective. $\square$
\end{proof}

\begin{lemma}[the exponential map of a complete manifold is $C^\infty$]
  \label{lem:dc-ch7-3-expmap-smooth}
  \dcref{ch7:3.1}
  \uses{lem:dc-ch7-3-expmap-global}
  \lean{Riemannian.Exponential.contMDiff_expMapIntrinsic,
    Riemannian.Exponential.contMDiff_expMapGlobal,
    Riemannian.Jacobi.exists_geodesic_flow_step_contDiff,
    Riemannian.Jacobi.exists_geodesic_flow_step_contDiff_manifold_ball,
    Riemannian.Jacobi.exists_geodesic_flowstep_partition_contDiff,
    Riemannian.Jacobi.exists_geodesic_contDiff_chain_nbhd,
    Riemannian.Jacobi.contDiffAt_stateTransition}
  \leanok
  Let $M$ be a complete Riemannian manifold and $p\in M$. Then the exponential map
  $\exp_p:T_pM\to M$ is $C^\infty$. This is the smooth dependence of the geodesic flow on its
  initial conditions, specialised to the exponential map.
\end{lemma}
\begin{proof}
  \leanok
  \uses{lem:dc-ch7-3-expmap-global}
  Fix $v\in T_pM$ and let $\gamma=\gamma_v$ be the geodesic issuing from $p$ with velocity
  $v$. The compact arc $\gamma([0,1])$ is covered by finitely many charts; choose a partition
  $0=t_0\le t_1\le\cdots\le t_N=1$ so that each piece $[t_i,t_{i+1}]$ lies in a single chart
  $\beta_i$, with $\gamma(t_i),\gamma(t_{i+1})$ both in that chart. On each piece the geodesic
  flow of the coordinate spray is $C^\infty$ in its initial condition even away from the zero
  section: bump-extending the chart spray to a globally $C^\infty$ field and confining the
  flow to a small ball where the two agree, the flow inherits the smoothness of that field's
  Picard integral operator. The time-$t_{i+1}$ endpoint map of that flow is thus $C^\infty$ at
  the chart state of $\gamma$ at $t_i$, and it carries the state of \emph{any} nearby geodesic
  to that geodesic's state at $t_{i+1}$. The chart junction $(x,w)\mapsto(\tau(x),(D\tau)(x)w)$
  between overlapping charts is $C^\infty$ (the chart transition $\tau$ is $C^\infty$, so its
  value and the fibre-application of its derivative are). Composing the $N$ flow steps and
  $N-1$ junctions gives a map $F_0$, $C^\infty$ at $\gamma$'s initial chart state, whose value
  on a neighbourhood is the chart-$\beta_N$ state at time $1$ of the geodesic passing through
  the given initial state.

  Feed the family $w\mapsto\gamma_w$ of geodesics issuing from $p$ with velocity $w$. Since
  $\gamma_w$ is, by construction, the geodesic whose chart-$p$ velocity at time $0$ is $w$, its
  initial chart-$\beta_0$ state is $\iota(w)=(\varphi_{\beta_0}(p),\;C\,w)$, where
  $C$ is the fixed tangent coordinate change from the chart at $p$ to the chart at $\beta_0$;
  thus $\iota$ is affine, hence $C^\infty$, and no differentiation of $\exp_p$ near $0$ is
  needed. On a neighbourhood of $v$ one therefore has
  $\varphi_{\beta_N}(\exp_p(w))=\big(F_0(\iota(w))\big)_1$, a composition of $C^\infty$ maps,
  so the chart reading of $\exp_p$ is $C^\infty$ near $v$; reading back through the terminal
  chart (whose source contains $\exp_p(w)$ for $w$ near $v$) shows $\exp_p$ is $C^\infty$ at
  $v$. As $v$ was arbitrary, $\exp_p$ is $C^\infty$. $\square$
\end{proof}

\begin{lemma}[Hadamard assembly from geodesic completeness at a point]
  \label{lem:dc-ch7-3-1-assembly}
  \dcref{ch7:3.1}
  \uses{lem:dc-ch7-3-1-covering-diffeo, lem:dc-ch7-2-8-a-to-b, lem:dc-ch7-3-3}
  \lean{Riemannian.DCExpandsMetric.diffeomorphOfSimplyConnectedOfGeodesicCompleteAt}
  \leanok
  Let $f:M\to N$ be a smooth local diffeomorphism between Riemannian manifolds which
  \emph{expands the metric} ($|df_p(v)|\ge|v|$ for all $p\in M$ and all $v\in T_pM$), let
  $N$ be simply connected, and suppose $M$ is connected and \emph{geodesically complete at
  some point $o$}: every $v\in T_oM$ generates a geodesic $\gamma:\mathbb R\to M$, defined
  for all time, with $\gamma(0)=o$ and $\gamma'(0)=v$. Then $f$ is a diffeomorphism.

  This is \cref{lem:dc-ch7-3-1-covering-diffeo} with the completeness hypothesis stated as
  do Carmo states it in the Hadamard proof: instead of assuming $M$ proper, one \emph{derives}
  properness from geodesic completeness at $o$ (\cref{lem:dc-ch7-2-8-a-to-b}). It is the form
  applied to $f=\exp_p:T_pM\to M$, where $o$ is the origin of $T_pM$ and the hypothesis is
  precisely that the rays through the origin are geodesics of the pulled-back metric.
\end{lemma}
\begin{proof}
  \leanok
  \uses{lem:dc-ch7-3-1-covering-diffeo, lem:dc-ch7-2-8-a-to-b}
  By \cref{lem:dc-ch7-2-8-a-to-b}, geodesic completeness of $M$ at $o$ makes the closed and
  bounded subsets of $M$ compact, so $M$ is a proper — hence complete — metric space. Now
  \cref{lem:dc-ch7-3-1-covering-diffeo} applies verbatim: the metric-expanding smooth local
  diffeomorphism $f$ out of the complete manifold $M$ onto the simply connected manifold $N$
  is a covering map (\cref{lem:dc-ch7-3-3}), and a covering map onto a simply connected base
  is a diffeomorphism. $\square$
\end{proof}

\subsection{Proof of the Hadamard theorem}
Since $M$ is complete, $\exp_p:T_pM\to M$ is defined for all $p\in M$ and is
surjective. By \cref{lem:dc-ch7-3-2}, $\exp_p$ is a local diffeomorphism. This allows us to
introduce a Riemannian metric on $T_pM$ in such a way that $\exp_p$ is a local
isometry. Such a metric is complete, because the geodesics of $T_pM$ passing
through the origin are straight lines (cf. \cref{thm:dc-ch7-2-8}, (a) $\Rightarrow$ (d)):
this is the geodesic-completeness-at-the-origin hypothesis of \cref{lem:dc-ch7-3-1-assembly}.
Thus $\exp_p$ is a metric-expanding (indeed metric-preserving) local diffeomorphism
out of the manifold $T_pM$ — geodesically complete at the origin — onto the simply connected
manifold $M$, so by \cref{lem:dc-ch7-3-1-assembly} (which derives the properness of $T_pM$
from that completeness via \cref{lem:dc-ch7-2-8-a-to-b} and then packages the covering-map
conclusion of \cref{lem:dc-ch7-3-3} with the simply-connectedness of $M$) it is a
diffeomorphism. The two ingredients still required are that $\exp_p$ is a local
diffeomorphism (\cref{lem:dc-ch7-3-2}, from the absence of conjugate points when $K\le 0$)
and that the rays of $T_pM$ are geodesics of the pulled-back metric; the entire
topological and smooth remainder is \cref{lem:dc-ch7-3-1-assembly}.
$\square$

\begin{lemma}[the tangent space with the pulled-back metric]
  \label{lem:dc-ch7-3-4-pullback-metric}
  \dcref{ch7:3.4}
  \uses{def:dc-ch7-2-4, def:dc-ch1-2-2, lem:dc-ch7-3-3-length-expansion}
  \lean{Riemannian.HadamardModel, Riemannian.HadamardModel.flatMetric, Riemannian.HadamardModel.pullbackMetric, Riemannian.HadamardModel.dcExpandsMetric_pullbackMetric}
  \leanok
  Let $M$ be a Riemannian manifold, $p\in M$, and $f:T_pM\to M$ a smooth local
  diffeomorphism (for instance $f=\exp_p$ at a pole). Regard the tangent space $T_pM$ as
  a smooth manifold in its own right, modelled on $\mathbb R^n=T_pM$ via the single global
  identity chart. Then $T_pM$ carries the \emph{pulled-back Riemannian metric} $f^*g$,
  $\langle u,v\rangle^{f^*g}_x=\langle df_x(u),df_x(v)\rangle_{f(x)}$, for which $f$ is a
  local isometry, hence a metric-expander: $|df_x(v)|=|v|$ for all $x,v$.
\end{lemma}
\begin{proof}
  \leanok
  \uses{lem:dc-ch7-3-3-length-expansion}
  The pulled-back form is a Riemannian metric on $T_pM$ because $f$ is a smooth
  immersion (its differential is injective, being a linear isomorphism at each point):
  symmetry and positivity are inherited from $g$ through the injective differential, and
  the metric depends smoothly on the base point since $f$ is smooth. By construction
  $\langle u,v\rangle^{f^*g}_x=\langle df_x u,df_x v\rangle_{f(x)}$, so $f$ preserves the
  metric and in particular $|df_x(v)|\ge|v|$ for all $v$
  (\cref{lem:dc-ch7-3-3-length-expansion}). $\square$
\end{proof}

\begin{lemma}[pole $\Rightarrow$ diffeomorphism]
  \label{lem:dc-ch7-3-4-pole-diffeo}
  \dcref{ch7:3.4}
  \uses{lem:dc-ch7-3-1-assembly, lem:dc-ch7-3-4-pullback-metric, lem:dc-ch7-3-expmap-global}
  \lean{Riemannian.HadamardModel.diffeomorphOfPole, Riemannian.HadamardModel.diffeomorphOfPole_coe, Riemannian.HadamardModel.expDiffeomorphOfPole}
  \leanok
  Let $M$ be a complete, simply connected Riemannian manifold, $p\in M$, and $f:T_pM\to M$
  a smooth local diffeomorphism which is metric-expanding for its own pulled-back metric
  $f^*g$. If the radial lines $s\mapsto sv$ of $T_pM$ are geodesics of $f^*g$, then $f$ is a
  diffeomorphism $T_pM\to M$. For $f=\exp_p$ this is do Carmo's poles statement modulo its
  two analytic inputs: that $\exp_p$ is a local diffeomorphism is the pole hypothesis, and
  that the radial lines are $f^*g$-geodesics is the local-isometry observation.
\end{lemma}
\begin{proof}
  \leanok
  \uses{lem:dc-ch7-3-1-assembly, lem:dc-ch7-3-4-pullback-metric}
  Equip $T_pM$ with the pulled-back metric $f^*g$
  (\cref{lem:dc-ch7-3-4-pullback-metric}), for which $f$ is a metric-expanding smooth local
  diffeomorphism onto the simply connected manifold $M$. The radial lines through the origin
  are geodesics of $f^*g$ by hypothesis, defined for all time, so $T_pM$ is geodesically
  complete at the origin. \Cref{lem:dc-ch7-3-1-assembly} then derives the properness of
  $T_pM$ from that completeness and makes $f$ a diffeomorphism. $\square$
\end{proof}

\begin{lemma}[the poles theorem reduces to the rays being pullback-geodesics]
  \label{lem:dc-ch7-3-4-ray-reduction}
  \dcref{ch7:3.4}
  \uses{lem:dc-ch7-3-4-pole-diffeo, lem:dc-ch7-3-4-pullback-metric}
  \lean{Riemannian.HadamardModel.diffeomorphOfPole_of_rayGeodesic,
    Riemannian.HadamardModel.diffeomorphOfPole_of_rayGeodesic_coe,
    Riemannian.HadamardModel.hrays_of_rayGeodesic,
    Riemannian.HadamardModel.isGeodesic_rayCurve_of_christoffel,
    Riemannian.HadamardModel.rayCurve}
  \leanok
  Let $M$ be a complete, simply connected Riemannian manifold, $p\in M$, and
  $f:T_pM\to M$ a smooth local diffeomorphism which is metric-expanding for its own
  pulled-back metric $f^*g$. If for every $v\in T_pM$ the radial line
  $\gamma_v(s)=s\,v$ is a geodesic of $f^*g$, then $f$ is a diffeomorphism
  $T_pM\to M$. This isolates the single geometric input of
  \cref{lem:dc-ch7-3-4-pole-diffeo}: the radial lines being $f^*g$-geodesics.
\end{lemma}
\begin{proof}
  \leanok
  \uses{lem:dc-ch7-3-4-pole-diffeo}
  Regard $T_pM$ as the single-chart manifold on which the extended chart is the
  identity of the model space; then the chart reading of $\gamma_v$ is $s\mapsto s\,v$,
  whose second-order term vanishes, so the intrinsic geodesic equation of $f^*g$ along
  $\gamma_v$ reduces to the pointwise Christoffel identity
  $\Gamma^{f^*g}_{sv}(v,v)=0$. The remaining three requirements of the ray hypothesis
  of \cref{lem:dc-ch7-3-4-pole-diffeo} — that $\gamma_v$ starts at the origin, has
  chart-velocity $v$ at $0$, and is continuous — are immediate for the straight line
  $\gamma_v$. Hence the pole diffeomorphism of \cref{lem:dc-ch7-3-4-pole-diffeo}
  applies as soon as the radial lines are $f^*g$-geodesics. $\square$
\end{proof}

\begin{lemma}[$\exp_p$ is a local isometry, so the radial lines are $\exp_p^*g$-geodesics]
  \label{lem:dc-ch7-3-4-rays-are-geodesics}
  \dcref{ch7:3.4}
  \uses{lem:dc-ch7-3-4-pullback-metric, def:dc-ch7-2-2, def:dc-ch1-2-2, lem:dc-ch7-3-expmap-global}
  \lean{Riemannian.HadamardModel.expMap_rays_are_geodesics}
  \leanok
  For $f=\exp_p$ and the pulled-back metric $\exp_p^*g$ on $T_pM$, every radial line
  $\gamma_v(s)=s\,v$ is a geodesic of $\exp_p^*g$.
\end{lemma}
\begin{proof}
  \uses{lem:dc-ch7-3-4-pullback-metric, lem:dc-ch7-3-expmap-global}
  \leanok
  By construction $\exp_p:(T_pM,\exp_p^*g)\to(M,g)$ is a \emph{local isometry}: its
  differential preserves the inner product (\cref{lem:dc-ch7-3-4-pullback-metric}).
  The image $\exp_p(\gamma_v(s))=\exp_p(s\,v)$ is, by definition of the exponential map,
  the geodesic of $M$ through $p$ with initial velocity $v$, defined for all $s$ since
  $M$ is complete. A local isometry \emph{reflects} geodesics: because $d(\exp_p)$
  is a linear isomorphism at every point that intertwines the Levi-Civita connections
  of $\exp_p^*g$ and $g$ (the connection is natural under isometry), a curve $\gamma$
  is a geodesic of $\exp_p^*g$ if and only if $\exp_p\circ\gamma$ is a geodesic of $g$.
  Applying this to $\gamma=\gamma_v$, whose image is the $g$-geodesic $s\mapsto\exp_p(s\,v)$,
  shows $\gamma_v$ is a geodesic of $\exp_p^*g$. $\square$
\end{proof}

\begin{remark}[poles]
  \label{rem:dc-ch7-3-4}
  \dcref{ch7:3.4}
  \uses{lem:dc-ch7-3-4-pole-diffeo, lem:dc-ch7-3-4-ray-reduction, lem:dc-ch7-3-4-rays-are-geodesics, lem:dc-ch7-3-expmap-global}
  \lean{Riemannian.HadamardModel.expDiffeomorphOfPole_of_pole}
  \leanok
  The above proof gives a little more than what is stated. Call a point $p$ of a
  complete Riemannian manifold $M$ a \emph{pole} if it has the property that it has no
  conjugate points. Any point of a complete manifold $M$ with non-positive
  sectional curvature is a pole of $M$. However, poles can exist in non-compact
  manifolds which have positive sectional curvature (see Exercise 13). What we have
  just proved is the following general fact. *If a complete simply connected
  Riemannian manifold $M$ has a pole, then $M$ is diffeomorphic to $\mathbb{R}^n$,
  $n=\dim M$.* The reduction of this statement to its two analytic inputs is
  \cref{lem:dc-ch7-3-4-pole-diffeo}: it remains only to supply, for $f=\exp_p$, that
  $\exp_p$ is a smooth local diffeomorphism (the pole hypothesis) and that the radial lines
  of $T_pM$ are geodesics of the pulled-back metric $\exp_p^*g$ (that $\exp_p$ is a local
  isometry reflects geodesics from $M$ to $T_pM$); the entire topological and smooth
  remainder is discharged. Formally, \cref{lem:dc-ch7-3-4-ray-reduction} discharges the
  three regularity inputs of the ray hypothesis and reduces the poles theorem to the
  single geometric fact of \cref{lem:dc-ch7-3-4-rays-are-geodesics} — that the radial
  lines are $\exp_p^*g$-geodesics — whose Lean residual is the naturality of the
  Levi-Civita connection under a local isometry (a curve is an $\exp_p^*g$-geodesic iff
  its $\exp_p$-image is a $g$-geodesic).
\end{remark}
